\section{Isbell--Mrówka spaces}

\import{function.tex}
\import{set/finite.tex}
\import{set/almost-disjoint.tex}
\import{set/mad-family.tex}
\import{set/finite-ordinal.tex}
\import{set/function-space.tex}
\import{set/choice.tex}
\import{cardinal/countability.tex}
\import{topology/topological-space.tex}
\import{topology/basis.tex}
\import{topology/continuity.tex}
\import{topology/subspace.tex}
\import{topology/countability-axioms.tex}
\import{topology/discrete.tex}
\import{topology/compactness.tex}
\import{topology/separation.tex}
\import{topology/clopen-basis.tex}
\import{topology/complete-regularity.tex}
\import{topology/local-compactness.tex}
\import{topology/countable-compactness.tex}
\import{topology/pseudocompactness.tex}
\import{surreal/embedding.tex}
\import{surreal/real.tex}
\import{surreal/absolute-value.tex}
\import{surreal/real-metric.tex}
\import{surreal/real-boundedness.tex}
\import{topology/metric.tex}
\import{topology/real-line.tex}

\begin{definition}\label{psi_carrier}
    $\psiCarrier{P}=\naturals\union P$.
\end{definition}

\begin{definition}\label{psi_neighbourhood}
    $\psiNeighbourhood{A}{F}=\{A\}\union(A\setminus F)$.
\end{definition}

\begin{proposition}\label{psi_carrier_naturals_disjoint_family}
    Suppose $P$ is an almost-disjoint family on
        $\naturals$.
    Then $\naturals\inter P=\emptyset$.
\end{proposition}
\begin{proof}
    Show that every element of $\naturals\inter P$ is an element
        of $\emptyset$.
    \begin{subproof}
        Fix $x\in\naturals\inter P$.
        $x\in\naturals$ and $x\in P$ by \obv{inter}.
        Hence $x$ is finite by \obv{natural_number_is_finite}.
        $x$ is infinite by
            \obv{almost_disjoint_members_infinite}.
        Contradiction.
    \end{subproof}
    Follows by \auto{subseteq_emptyset_iff,subseteq}.
\end{proof}

\begin{proposition}\label{naturals_subset_psi_carrier}
    $\naturals\subseteq\psiCarrier{P}$.
\end{proposition}
\begin{proof}
    Follows by \auto{psi_carrier,union_upper_left}.
\end{proof}

\begin{proposition}\label{psi_family_subset_psi_carrier}
    $P\subseteq\psiCarrier{P}$.
\end{proposition}
\begin{proof}
    Follows by \auto{psi_carrier,union_upper_right}.
\end{proof}

\begin{proposition}\label{psi_carrier_iff}
    $x\in\psiCarrier{P}$ iff
        $x\in\naturals$ or $x\in P$.
\end{proposition}
\begin{proof}
    Follows by \auto{psi_carrier,union_iff}.
\end{proof}

\begin{proposition}\label{psi_carrier_point_cases}
    Suppose $P$ is an almost-disjoint family on
        $\naturals$.
    Let $x\in\psiCarrier{P}$.
    Then $x\in\naturals$ or $x\in P$.
\end{proposition}
\begin{proof}
    Follows by \obv{psi_carrier_iff}.
\end{proof}

\begin{proposition}\label{psi_neighbourhood_subset_carrier}
    Suppose $P$ is an almost-disjoint family on
        $\naturals$.
    Suppose $A\in P$.
    Then $\psiNeighbourhood{A}{F}\subseteq\psiCarrier{P}$.
\end{proposition}
\begin{proof}
    $P\subseteq\psiCarrier{P}$ by \obv{psi_family_subset_psi_carrier}.
    Hence $A\in\psiCarrier{P}$ by \obv{subseteq}.
    $A\subseteq\naturals$
        by \obv{almost_disjoint_members_subset}.
    Hence $A\setminus F\subseteq\naturals$
        by \obv{setminus_subseteq,subseteq_transitive}.
    Hence $A\setminus F\subseteq\psiCarrier{P}$
        by \auto{naturals_subset_psi_carrier,subseteq_transitive}.
    Follows by
        \auto{psi_neighbourhood,union_iff,singleton_iff,subseteq}.
\end{proof}

\begin{proposition}\label{psi_neighbourhood_contains_center}
    $A\in\psiNeighbourhood{A}{F}$.
\end{proposition}
\begin{proof}
    Follows by
        \obv{psi_neighbourhood,union_intro_left,singleton_intro}.
\end{proof}

\begin{proposition}\label{psi_neighbourhood_natural_trace}
    Suppose $A\subseteq\naturals$ and $A$ is infinite.
    Then $\psiNeighbourhood{A}{F}\inter\naturals=A\setminus F$.
\end{proposition}
\begin{proof}
    Show that $A\notin\naturals$.
    \begin{subproof}
        Suppose $A\in\naturals$.
        Hence $A$ is finite by \obv{natural_number_is_finite}.
        Contradiction.
    \end{subproof}
    Show that every element of
        $\psiNeighbourhood{A}{F}\inter\naturals$
        is an element of $A\setminus F$.
    \begin{subproof}
        Fix $x\in\psiNeighbourhood{A}{F}\inter\naturals$.
        $x\in\psiNeighbourhood{A}{F}$ and $x\in\naturals$
            by \obv{inter}.
        $x=A$ or $x\in A\setminus F$
            by \auto{psi_neighbourhood,union_iff,singleton_iff}.
        $x\neq A$.
    \end{subproof}
    Hence $\psiNeighbourhood{A}{F}\inter\naturals
        \subseteq A\setminus F$ by \obv{subseteq}.
    Show that every element of $A\setminus F$ is an element of
        $\psiNeighbourhood{A}{F}\inter\naturals$.
    \begin{subproof}
        Fix $x\in A\setminus F$.
        $x\in A$ by \obv{setminus}.
        Hence $x\in\naturals$ by \obv{subseteq}.
        $x\in\psiNeighbourhood{A}{F}$
            by \obv{psi_neighbourhood,union_intro_right}.
        Follows by \obv{inter}.
    \end{subproof}
    Hence $A\setminus F\subseteq
        \psiNeighbourhood{A}{F}\inter\naturals$ by \obv{subseteq}.
    Follows by \obv{subseteq_antisymmetric}.
\end{proof}

\begin{proposition}\label{psi_neighbourhood_family_trace}
    Suppose $P$ is an almost-disjoint family on
        $\naturals$.
    Suppose $A\in P$.
    Then $\psiNeighbourhood{A}{F}\inter P=\{A\}$.
\end{proposition}
\begin{proof}
    $A\subseteq\naturals$
        by \obv{almost_disjoint_members_subset}.
    $A\setminus F\subseteq\naturals$
        by \obv{setminus_subseteq,subseteq_transitive}.
    $\naturals\inter P=\emptyset$
        by \obv{psi_carrier_naturals_disjoint_family}.
    Follows by
        \auto{psi_neighbourhood,setext,inter,union_iff,singleton_iff,
            subseteq,emptyset}.
\end{proof}

\begin{proposition}\label{psi_neighbourhood_antitone}
    Suppose $F\subseteq G$.
    Then $\psiNeighbourhood{A}{G}\subseteq
        \psiNeighbourhood{A}{F}$.
\end{proposition}
\begin{proof}
    Follows by
        \auto{psi_neighbourhood,setminus,union_iff,singleton_iff,subseteq}.
\end{proof}

\begin{proposition}\label{psi_neighbourhood_same_center_intersection}
    $\psiNeighbourhood{A}{F}\inter\psiNeighbourhood{A}{G}
        =\psiNeighbourhood{A}{F\union G}$.
\end{proposition}
\begin{proof}
    Follows by
        \auto{psi_neighbourhood,setext,inter,union_iff,singleton_iff,
            setminus}.
\end{proof}

\begin{proposition}\label{psi_neighbourhood_distinct_centers_intersection}
    Suppose $P$ is an almost-disjoint family on
        $\naturals$.
    Suppose $A,B\in P$ and $A\neq B$.
    Then $\psiNeighbourhood{A}{F}\inter
        \psiNeighbourhood{B}{G}\subseteq A\inter B$.
\end{proposition}
\begin{proof}
    Show that every element of
        $\psiNeighbourhood{A}{F}\inter\psiNeighbourhood{B}{G}$
        is an element of $A\inter B$.
    \begin{subproof}
        Fix $x\in
            \psiNeighbourhood{A}{F}\inter\psiNeighbourhood{B}{G}$.
        $x\in\psiNeighbourhood{A}{F}$ and
            $x\in\psiNeighbourhood{B}{G}$ by \obv{inter}.
        $x=A$ or $x\in A\setminus F$
            by \auto{psi_neighbourhood,union_iff,singleton_iff}.
        $x=B$ or $x\in B\setminus G$
            by \auto{psi_neighbourhood,union_iff,singleton_iff}.
        Show that $x\neq A$.
        \begin{subproof}
            Suppose $x=A$.
            $A\in\psiNeighbourhood{B}{G}\inter P$
                by \obv{inter}.
            $\psiNeighbourhood{B}{G}\inter P=\{B\}$
                by \obv{psi_neighbourhood_family_trace}.
            Hence $A=B$ by \obv{singleton_iff}.
            Contradiction.
        \end{subproof}
        Show that $x\neq B$.
        \begin{subproof}
            Suppose $x=B$.
            $B\in\psiNeighbourhood{A}{F}\inter P$
                by \obv{inter}.
            $\psiNeighbourhood{A}{F}\inter P=\{A\}$
                by \obv{psi_neighbourhood_family_trace}.
            Hence $B=A$ by \obv{singleton_iff}.
            Contradiction.
        \end{subproof}
        Hence $x\in A\setminus F$ and $x\in B\setminus G$.
        Follows by \obv{setminus,inter}.
    \end{subproof}
    Follows by \obv{subseteq}.
\end{proof}

\begin{corollary}\label{psi_neighbourhood_distinct_centers_finite_intersection}
    Suppose $P$ is an almost-disjoint family on
        $\naturals$.
    Suppose $A,B\in P$ and $A\neq B$.
    Then $\psiNeighbourhood{A}{F}\inter
        \psiNeighbourhood{B}{G}$ is finite.
\end{corollary}
\begin{proof}
    $\psiNeighbourhood{A}{F}\inter\psiNeighbourhood{B}{G}
        \subseteq A\inter B$
        by \obv{psi_neighbourhood_distinct_centers_intersection}.
    $A\inter B$ is finite by
        \obv{almost_disjoint_distinct_intersection_finite}.
    Follows by \obv{subseteq_finite_is_finite}.
\end{proof}

\begin{proposition}\label{psi_neighbourhood_normalize_finite_set}
    $\psiNeighbourhood{A}{F}
        =\psiNeighbourhood{A}{F\inter A}$.
\end{proposition}
\begin{proof}
    Follows by
        \auto{psi_neighbourhood,setext,union_iff,singleton_iff,setminus,
            inter}.
\end{proof}

\begin{proposition}\label{psi_neighbourhood_avoids_natural_set}
    Suppose $A,B\subseteq\naturals$ and $A$ is infinite.
    Then $\psiNeighbourhood{A}{A\inter B}\inter B=\emptyset$.
\end{proposition}
\begin{proof}
    Show that $A\notin\naturals$.
    \begin{subproof}
        Suppose $A\in\naturals$.
        Hence $A$ is finite by \obv{natural_number_is_finite}.
        Contradiction.
    \end{subproof}
    Follows by
        \auto{psi_neighbourhood,setext,union_iff,singleton_iff,setminus,
            inter,subseteq,emptyset}.
\end{proof}

\begin{definition}\label{psi_natural_basis}
    $\psiNaturalBasis=\{\{n\}\mid n\in\naturals\}$.
\end{definition}

\begin{definition}\label{psi_family_basis}
    $\psiFamilyBasis{P}=\{\psiNeighbourhood{A}{F}
        \mid A\in P,F\in\finiteSubsets{\naturals}\}$.
\end{definition}

\begin{definition}\label{psi_basis}
    $\psiBasis{P}=\psiNaturalBasis\union\psiFamilyBasis{P}$.
\end{definition}

\begin{proposition}\label{psi_natural_singleton_in_basis}
    Let $n\in\naturals$.
    Then $\{n\}\in\psiBasis{P}$.
\end{proposition}
\begin{proof}
    $\{n\}\in\psiNaturalBasis$ by \obv{psi_natural_basis}.
    Follows by \obv{psi_basis,union_intro_left}.
\end{proof}

\begin{proposition}\label{psi_neighbourhood_in_basis}
    Suppose $A\in P$ and $F\subseteq\naturals$ and $F$ is finite.
    Then $\psiNeighbourhood{A}{F}\in\psiBasis{P}$.
\end{proposition}
\begin{proof}
    $F\in\finiteSubsets{\naturals}$ by \auto{finite_subsets_iff}.
    Hence $\psiNeighbourhood{A}{F}\in\psiFamilyBasis{P}$
        by \obv{psi_family_basis}.
    Follows by \obv{psi_basis,union_intro_right}.
\end{proof}

\begin{proposition}\label{psi_basis_member_cases}
    Let $V\in\psiBasis{P}$.
    Then (there exists $n\in\naturals$ such that $V=\{n\}$) or
        (there exist $A,F$ such that $A\in P$ and
            $F\in\finiteSubsets{\naturals}$ and
            $V=\psiNeighbourhood{A}{F}$).
\end{proposition}
\begin{proof}
    $V\in\psiNaturalBasis$ or $V\in\psiFamilyBasis{P}$
        by \auto{psi_basis,union_iff}.
    Follows by \auto{psi_natural_basis,psi_family_basis}.
\end{proof}

\begin{proposition}\label{psi_basis_member_subset_carrier}
    Suppose $P$ is an almost-disjoint family on $\naturals$.
    Let $V\in\psiBasis{P}$.
    Then $V\subseteq\psiCarrier{P}$.
\end{proposition}
\begin{proof}
    (There exists $n\in\naturals$ such that $V=\{n\}$) or
        (there exist $A,F$ such that $A\in P$ and
            $F\in\finiteSubsets{\naturals}$ and
            $V=\psiNeighbourhood{A}{F}$)
        by \obv{psi_basis_member_cases}.
    \begin{byCase}
        \caseOf{There exists $n\in\naturals$ such that $V=\{n\}$.}
            Take $n\in\naturals$ such that $V=\{n\}$.
            $n\in\psiCarrier{P}$ by
                \auto{naturals_subset_psi_carrier,subseteq}.
            Hence $\{n\}\subseteq\psiCarrier{P}$
                by \obv{singleton_subset_intro}.
        \caseOf{There exist $A,F$ such that $A\in P$ and
                $F\in\finiteSubsets{\naturals}$ and
                $V=\psiNeighbourhood{A}{F}$.}
            Take $A,F$ such that $A\in P$ and
                $F\in\finiteSubsets{\naturals}$ and
                $V=\psiNeighbourhood{A}{F}$.
            Follows by \auto{psi_neighbourhood_subset_carrier}.
    \end{byCase}
\end{proof}

\begin{proposition}\label{psi_basis_covers_carrier}
    Suppose $P$ is an almost-disjoint family on $\naturals$.
    Then $\psiBasis{P}$ covers $\psiCarrier{P}$.
\end{proposition}
\begin{proof}
    Show that for every $x\in\psiCarrier{P}$ there exists
        $V\in\psiBasis{P}$ such that $x\in V$.
    \begin{subproof}
        Fix $x\in\psiCarrier{P}$.
        $x\in\naturals$ or $x\in P$ by \obv{psi_carrier_point_cases}.
        \begin{byCase}
            \caseOf{$x\in\naturals$.}
                $\{x\}\in\psiBasis{P}$
                    by \obv{psi_natural_singleton_in_basis}.
                $x\in\{x\}$ by \obv{singleton_intro}.
            \caseOf{$x\in P$.}
                $\emptyset\subseteq\naturals$ by \obv{emptyset_subseteq}.
                $\emptyset$ is finite by \obv{emptyset_is_finite}.
                Hence $\psiNeighbourhood{x}{\emptyset}\in\psiBasis{P}$
                    by \obv{psi_neighbourhood_in_basis}.
                $x\in\psiNeighbourhood{x}{\emptyset}$
                    by \obv{psi_neighbourhood_contains_center}.
        \end{byCase}
    \end{subproof}
\end{proof}

\begin{proposition}\label{psi_family_point_in_basis_member}
    Suppose $P$ is an almost-disjoint family on $\naturals$.
    Suppose $V\in\psiBasis{P}$ and $A\in P\inter V$.
    Then there exists $F$ such that
        $F\subseteq\naturals$ and $F$ is finite and
        $V=\psiNeighbourhood{A}{F}$.
\end{proposition}
\begin{proof}
    (There exists $n\in\naturals$ such that $V=\{n\}$) or
        (there exist $B,F$ such that $B\in P$ and
            $F\in\finiteSubsets{\naturals}$ and
            $V=\psiNeighbourhood{B}{F}$)
        by \obv{psi_basis_member_cases}.
    \begin{byCase}
        \caseOf{There exists $n\in\naturals$ such that $V=\{n\}$.}
            Take $n\in\naturals$ such that $V=\{n\}$.
            $A\in P$ and $A\in V$ by \obv{inter}.
            Hence $A=n$ by \obv{singleton_iff}.
            Hence $A\in\naturals\inter P$ by \obv{inter}.
            $\naturals\inter P=\emptyset$
                by \obv{psi_carrier_naturals_disjoint_family}.
            Contradiction by \obv{emptyset}.
        \caseOf{There exist $B,F$ such that $B\in P$ and
                $F\in\finiteSubsets{\naturals}$ and
                $V=\psiNeighbourhood{B}{F}$.}
            Take $B,F$ such that $B\in P$ and
                $F\in\finiteSubsets{\naturals}$ and
                $V=\psiNeighbourhood{B}{F}$.
            $A\in P$ and $A\in V$ by \obv{inter}.
            Hence $A\in\psiNeighbourhood{B}{F}\inter P$
                by \obv{inter}.
            $\psiNeighbourhood{B}{F}\inter P=\{B\}$
                by \obv{psi_neighbourhood_family_trace}.
            Hence $A=B$ by \obv{singleton_iff}.
            $F\subseteq\naturals$ and $F$ is finite
                by \obv{finite_subsets_iff}.
            Follows by \auto.
    \end{byCase}
\end{proof}

\begin{theorem}\label{psi_basis_is_topological_basis}
    Suppose $P$ is an almost-disjoint family on $\naturals$.
    Then $\psiBasis{P}$ is a topological basis for $\psiCarrier{P}$.
\end{theorem}
\begin{proof}
    Show that $\psiBasis{P}$ is a family of subsets of $\psiCarrier{P}$.
    \begin{subproof}
        Show that every element of $\psiBasis{P}$ is an element of
            $\pow{\psiCarrier{P}}$.
        \begin{subproof}
            Fix $V\in\psiBasis{P}$.
            $V\subseteq\psiCarrier{P}$
                by \obv{psi_basis_member_subset_carrier}.
            Follows by \obv{pow_iff,subseteq}.
        \end{subproof}
        Follows by \auto{subseteq}.
    \end{subproof}
    $\psiBasis{P}$ covers $\psiCarrier{P}$
        by \obv{psi_basis_covers_carrier}.
    Show that for all $U,V,x$ such that
        $U,V\in\psiBasis{P}$ and $x\in U,V$
        there exists $W\in\psiBasis{P}$ such that
            $x\in W\subseteq U,V$.
    \begin{subproof}
        Fix $U,V,x$.
        Assume $U,V\in\psiBasis{P}$ and $x\in U,V$.
        $U,V\subseteq\psiCarrier{P}$
            by \obv{psi_basis_member_subset_carrier}.
        Hence $x\in\psiCarrier{P}$ by \obv{subseteq}.
        $x\in\naturals$ or $x\in P$ by \obv{psi_carrier_point_cases}.
        \begin{byCase}
            \caseOf{$x\in\naturals$.}
                $\{x\}\in\psiBasis{P}$
                    by \obv{psi_natural_singleton_in_basis}.
                $x\in\{x\}$ by \obv{singleton_intro}.
                $\{x\}\subseteq U,V$ by \obv{singleton_subset_intro}.
            \caseOf{$x\in P$.}
                $x\in P\inter U$ and $x\in P\inter V$
                    by \obv{inter}.
                Take $F$ such that $F\subseteq\naturals$ and $F$ is finite
                    and $U=\psiNeighbourhood{x}{F}$
                    by \obv{psi_family_point_in_basis_member}.
                Take $G$ such that $G\subseteq\naturals$ and $G$ is finite
                    and $V=\psiNeighbourhood{x}{G}$
                    by \obv{psi_family_point_in_basis_member}.
                $F\union G\subseteq\naturals$
                    by \auto{union_iff,subseteq}.
                $F\union G$ is finite by \obv{union_of_finite_is_finite}.
                Hence $\psiNeighbourhood{x}{F\union G}\in\psiBasis{P}$
                    by \obv{psi_neighbourhood_in_basis}.
                $x\in\psiNeighbourhood{x}{F\union G}$
                    by \obv{psi_neighbourhood_contains_center}.
                $\psiNeighbourhood{x}{F\union G}=U\inter V$
                    by \auto{psi_neighbourhood_same_center_intersection}.
                Hence $\psiNeighbourhood{x}{F\union G}\subseteq U,V$
                    by \auto{inter,subseteq}.
        \end{byCase}
    \end{subproof}
    Follows by \obv{topological_basis}.
\end{proof}

\begin{definition}\label{psi_space}
    $\psiSpace{P}=\topologyGeneratedByBasis{\psiBasis{P}}{\psiCarrier{P}}$.
\end{definition}

\begin{classification}\label{psi_space_is_topological_space}
    Suppose $P$ is an almost-disjoint family on $\naturals$.
    Then $\psiSpace{P}$ is a topological space.
\end{classification}
\begin{proof}
    $\psiBasis{P}$ is a topological basis for $\psiCarrier{P}$
        by \obv{psi_basis_is_topological_basis}.
    Follows by
        \auto{psi_space,topology_generated_by_basis_is_topological_space}.
\end{proof}

\begin{proposition}\label{psi_space_carrier}
    $\carrier[\psiSpace{P}]=\psiCarrier{P}$.
\end{proposition}
\begin{proof}
    Follows by
        \auto{psi_space,topology_generated_by_basis_carrier}.
\end{proof}

\begin{proposition}\label{psi_space_opens}
    Suppose $P$ is an almost-disjoint family on $\naturals$.
    $\opens[\psiSpace{P}]=\genOpens{\psiBasis{P}}{\psiCarrier{P}}$.
\end{proposition}
\begin{proof}
    $\psiBasis{P}$ is a topological basis for $\psiCarrier{P}$ by
        \obv{psi_basis_is_topological_basis}.
    Follows by
        \auto{psi_space,topology_generated_by_basis_opens}.
\end{proof}

\begin{theorem}\label{psi_basis_is_open_basis}
    Suppose $P$ is an almost-disjoint family on $\naturals$.
    Then $\psiBasis{P}$ is an open basis of $\psiSpace{P}$.
\end{theorem}
\begin{proof}
    $\psiBasis{P}$ is a topological basis for $\psiCarrier{P}$
        by \obv{psi_basis_is_topological_basis}.
    Follows by
        \auto{psi_space,topology_generated_by_basis_has_open_basis}.
\end{proof}

\begin{theorem}\label{psi_open_iff}
    Suppose $P$ is an almost-disjoint family on $\naturals$.
    Then $U$ is open in $\psiSpace{P}$ iff
        $U\subseteq\psiCarrier{P}$ and
        for every $A\in P\inter U$ we have $A\setminus U$ is finite.
\end{theorem}
\begin{proof}
    Show that if $U$ is open in $\psiSpace{P}$, then
        $U\subseteq\psiCarrier{P}$ and
        for every $A\in P\inter U$ we have $A\setminus U$ is finite.
    \begin{subproof}
        Assume $U$ is open in $\psiSpace{P}$.
        $U\in\opens[\psiSpace{P}]$ by \obv{is_open}.
        Hence $U\subseteq\carrier[\psiSpace{P}]$ by \obv{opens_type}.
        Hence $U\subseteq\psiCarrier{P}$ by \obv{psi_space_carrier}.
        Show that for every $A\in P\inter U$ we have
            $A\setminus U$ is finite.
        \begin{subproof}
            Fix $A\in P\inter U$.
            $A\in U$ by \obv{inter}.
            $U\in\opens[\psiSpace{P}]$ by \obv{is_open}.
            $\psiBasis{P}$ is an open basis of $\psiSpace{P}$
                by \obv{psi_basis_is_open_basis}.
            Take $V\in\psiBasis{P}$ such that $A\in V\subseteq U$
                by \obv{open_basis}.
            $A\in P\inter V$ by \obv{inter}.
            Take $F$ such that $F\subseteq\naturals$ and $F$ is finite
                and $V=\psiNeighbourhood{A}{F}$
                by \obv{psi_family_point_in_basis_member}.
            Show that every element of $A\setminus U$ is an element of $F$.
            \begin{subproof}
                Fix $x\in A\setminus U$.
                Suppose not.
                Hence $x\notin F$.
                Hence $x\in A\setminus F$ by \obv{setminus}.
                Hence $x\in\psiNeighbourhood{A}{F}$
                    by \obv{psi_neighbourhood,union_intro_right}.
                Hence $x\in U$ by \obv{subseteq}.
                Contradiction by \obv{setminus}.
            \end{subproof}
            Hence $A\setminus U\subseteq F$ by \obv{subseteq}.
            Follows by \obv{subseteq_finite_is_finite}.
        \end{subproof}
    \end{subproof}
    Show that if $U\subseteq\psiCarrier{P}$ and
        for every $A\in P\inter U$ we have $A\setminus U$ is finite,
        then $U$ is open in $\psiSpace{P}$.
    \begin{subproof}
        Assume $U\subseteq\psiCarrier{P}$ and
            for every $A\in P\inter U$ we have $A\setminus U$ is finite.
        $U\in\pow{\psiCarrier{P}}$ by \obv{pow_iff,subseteq}.
        Show that for every $x\in U$ there exists $V\in\psiBasis{P}$
            such that $x\in V\subseteq U$.
        \begin{subproof}
            Fix $x\in U$.
            Hence $x\in\psiCarrier{P}$ by \obv{subseteq}.
            $x\in\naturals$ or $x\in P$
                by \obv{psi_carrier_point_cases}.
            \begin{byCase}
                \caseOf{$x\in\naturals$.}
                    $\{x\}\in\psiBasis{P}$
                        by \obv{psi_natural_singleton_in_basis}.
                    $x\in\{x\}$ by \obv{singleton_intro}.
                    $\{x\}\subseteq U$ by \obv{singleton_subset_intro}.
                \caseOf{$x\in P$.}
                    $x\in P\inter U$ by \obv{inter}.
                    Hence $x\setminus U$ is finite.
                    $x\subseteq\naturals$ by
                        \obv{almost_disjoint_members_subset}.
                    Hence $x\setminus U\subseteq\naturals$
                        by \obv{setminus_subseteq,subseteq_transitive}.
                    $\psiNeighbourhood{x}{x\setminus U}\in\psiBasis{P}$
                        by \obv{psi_neighbourhood_in_basis}.
                    $x\in\psiNeighbourhood{x}{x\setminus U}$
                        by \obv{psi_neighbourhood_contains_center}.
                    Show that every element of
                        $\psiNeighbourhood{x}{x\setminus U}$
                        is an element of $U$.
                    \begin{subproof}
                        Fix $y\in\psiNeighbourhood{x}{x\setminus U}$.
                        $y=x$ or $y\in x\setminus(x\setminus U)$
                            by
                            \auto{psi_neighbourhood,union_iff,
                                singleton_iff}.
                        \begin{byCase}
                            \caseOf{$y=x$.}
                                Follows.
                            \caseOf{$y\in x\setminus(x\setminus U)$.}
                                Follows by \auto{setminus}.
                        \end{byCase}
                    \end{subproof}
                    Hence $\psiNeighbourhood{x}{x\setminus U}\subseteq U$
                        by \obv{subseteq}.
            \end{byCase}
        \end{subproof}
        Hence $U\in\genOpens{\psiBasis{P}}{\psiCarrier{P}}$
            by \obv{genopens}.
        Hence $U\in\opens[\psiSpace{P}]$ by \obv{psi_space_opens}.
        $\psiSpace{P}$ is a topological space
            by \obv{psi_space_is_topological_space}.
        Follows by \obv{is_open}.
    \end{subproof}
    Follows by \auto.
\end{proof}

\begin{corollary}\label{psi_family_point_open_refinement}
    Suppose $P$ is an almost-disjoint family on $\naturals$.
    Suppose $U$ is open in $\psiSpace{P}$.
    Suppose $A\in P\inter U$.
    Then there exists $F$ such that
        $F\subseteq\naturals$ and $F$ is finite and
        $\psiNeighbourhood{A}{F}\subseteq U$.
\end{corollary}
\begin{proof}
    $U\in\opens[\psiSpace{P}]$ by \obv{is_open}.
    $A\in U$ by \obv{inter}.
    $\psiBasis{P}$ is an open basis of $\psiSpace{P}$
        by \obv{psi_basis_is_open_basis}.
    Take $V\in\psiBasis{P}$ such that $A\in V\subseteq U$
        by \obv{open_basis}.
    $A\in P\inter V$ by \obv{inter}.
    Take $F$ such that $F\subseteq\naturals$ and $F$ is finite and
        $V=\psiNeighbourhood{A}{F}$
        by \obv{psi_family_point_in_basis_member}.
    Follows by \auto.
\end{proof}

\begin{corollary}\label{psi_family_point_canonical_open_refinement}
    Suppose $P$ is an almost-disjoint family on $\naturals$.
    Suppose $U$ is open in $\psiSpace{P}$.
    Suppose $A\in P\inter U$.
    Then $A\setminus U$ is finite and
        $\psiNeighbourhood{A}{A\setminus U}\subseteq U$.
\end{corollary}
\begin{proof}
    $A\setminus U$ is finite by \auto{psi_open_iff}.
    Show that every element of $\psiNeighbourhood{A}{A\setminus U}$
        is an element of $U$.
    \begin{subproof}
        Fix $x\in\psiNeighbourhood{A}{A\setminus U}$.
        $x=A$ or $x\in A\setminus(A\setminus U)$
            by \auto{psi_neighbourhood,union_iff,singleton_iff}.
        \begin{byCase}
            \caseOf{$x=A$.}
                Follows by \auto{inter}.
            \caseOf{$x\in A\setminus(A\setminus U)$.}
                Follows by \auto{setminus}.
        \end{byCase}
    \end{subproof}
    Hence $\psiNeighbourhood{A}{A\setminus U}\subseteq U$
        by \obv{subseteq}.
\end{proof}

\begin{theorem}\label{psi_closed_iff}
    Suppose $P$ is an almost-disjoint family on $\naturals$.
    Then $C$ is closed in $\psiSpace{P}$ iff
        $C\subseteq\psiCarrier{P}$ and
        for every $A\in P\setminus C$ we have $A\inter C$ is finite.
\end{theorem}
\begin{proof}
    $\psiSpace{P}$ is a topological space
        by \obv{psi_space_is_topological_space}.
    Show that if $C$ is closed in $\psiSpace{P}$, then
        $C\subseteq\psiCarrier{P}$ and
        for every $A\in P\setminus C$ we have $A\inter C$ is finite.
    \begin{subproof}
        Assume $C$ is closed in $\psiSpace{P}$.
        $C\subseteq\carrier[\psiSpace{P}]$
            by \obv{closed_in_subseteq_carrier}.
        Hence $C\subseteq\psiCarrier{P}$ by \auto{psi_space_carrier}.
        Let $U=\psiCarrier{P}\setminus C$.
        $U=\carrier[\psiSpace{P}]\setminus C$
            by \auto{psi_space_carrier}.
        Hence $U$ is open in $\psiSpace{P}$
            by \auto{complement_of_closed_is_open}.
        Show that for every $A\in P\setminus C$ we have
            $A\inter C$ is finite.
        \begin{subproof}
            Fix $A\in P\setminus C$.
            $A\in P$ and $A\notin C$ by \obv{setminus}.
            $A\in\psiCarrier{P}$
                by \auto{psi_family_subset_psi_carrier,subseteq}.
            Hence $A\in U$ by \obv{setminus}.
            Hence $A\in P\inter U$ by \obv{inter}.
            Hence $A\setminus U$ is finite by \auto{psi_open_iff}.
            Show that $A\setminus U=A\inter C$.
            \begin{subproof}
                $A\subseteq\naturals$ by
                    \obv{almost_disjoint_members_subset}.
                Hence $A\subseteq\psiCarrier{P}$ by
                    \auto{naturals_subset_psi_carrier,subseteq_transitive}.
                $U\subseteq\psiCarrier{P}$ by \auto{setminus_subseteq}.
                Hence $A\setminus U=
                    A\inter(\psiCarrier{P}\setminus U)$
                    by \obv{setminus_eq_inter_complement}.
                $\psiCarrier{P}\setminus U=C$
                    by \auto{double_relative_complement}.
            \end{subproof}
        \end{subproof}
    \end{subproof}
    Show that if $C\subseteq\psiCarrier{P}$ and
        for every $A\in P\setminus C$ we have $A\inter C$ is finite,
        then $C$ is closed in $\psiSpace{P}$.
    \begin{subproof}
        Assume $C\subseteq\psiCarrier{P}$ and
            for every $A\in P\setminus C$ we have $A\inter C$ is finite.
        Let $U=\psiCarrier{P}\setminus C$.
        Show that $U$ is open in $\psiSpace{P}$.
        \begin{subproof}
            $U\subseteq\psiCarrier{P}$ by \auto{setminus_subseteq}.
            Show that for every $A\in P\inter U$ we have
                $A\setminus U$ is finite.
            \begin{subproof}
                Fix $A\in P\inter U$.
                $A\in P$ and $A\in U$ by \obv{inter}.
                Hence $A\notin C$ by \auto{setminus}.
                Hence $A\in P\setminus C$ by \obv{setminus}.
                Hence $A\inter C$ is finite.
                Show that $A\setminus U=A\inter C$.
                \begin{subproof}
                    $A\subseteq\naturals$ by
                        \obv{almost_disjoint_members_subset}.
                    Hence $A\subseteq\psiCarrier{P}$ by
                        \auto{naturals_subset_psi_carrier,
                            subseteq_transitive}.
                    $U\subseteq\psiCarrier{P}$
                        by \obv{setminus_subseteq}.
                    Hence $A\setminus U=
                        A\inter(\psiCarrier{P}\setminus U)$
                        by \obv{setminus_eq_inter_complement}.
                    $\psiCarrier{P}\setminus U=C$
                        by \auto{double_relative_complement}.
                \end{subproof}
            \end{subproof}
            Follows by \auto{psi_open_iff}.
        \end{subproof}
        $U$ is open in $\psiSpace{P}$.
        $U=\carrier[\psiSpace{P}]\setminus C$
            by \auto{psi_space_carrier}.
        Hence $U\subseteq\carrier[\psiSpace{P}]$
            by \auto{setminus_subseteq}.
        Hence $\carrier[\psiSpace{P}]\setminus U$ is closed in
            $\psiSpace{P}$ by \obv{complement_of_open_is_closed}.
        $\carrier[\psiSpace{P}]\setminus U=C$
            by
            \auto{psi_space_carrier,double_relative_complement}.
    \end{subproof}
    Follows by \auto.
\end{proof}

\begin{theorem}\label{psi_closure_formula}
    Suppose $P$ is an almost-disjoint family on $\naturals$.
    Then
        $\closure{C}{\psiSpace{P}}=
        (C\inter\psiCarrier{P})\union
        \{A\in P\mid \text{$A\inter C$ is infinite}\}$.
\end{theorem}
\begin{proof}
    Let $H=\{A\in P\mid \text{$A\inter C$ is infinite}\}$.
    Let $K=(C\inter\psiCarrier{P})\union H$.
    $H\subseteq P$ by \obv{subseteq}.
    Hence $H\subseteq\psiCarrier{P}$ by
        \auto{psi_family_subset_psi_carrier,subseteq_transitive}.
    $C\inter\psiCarrier{P}\subseteq\psiCarrier{P}$
        by \obv{inter_lower_right}.
    Hence $K\subseteq\psiCarrier{P}$
        by \auto{union_iff,subseteq}.
    $C\inter\psiCarrier{P}\subseteq K$ by \auto{union_upper_left}.
    Show that $K$ is closed in $\psiSpace{P}$.
    \begin{subproof}
        Show that for every $A\in P\setminus K$ we have
            $A\inter K$ is finite.
        \begin{subproof}
            Fix $A\in P\setminus K$.
            $A\in P$ and $A\notin K$ by \obv{setminus}.
            $A\notin H$ by \auto{union_iff}.
            Hence $A\inter C$ is finite by \auto.
            Show that every element of $A\inter K$ is an element of
                $A\inter C$.
            \begin{subproof}
                Fix $x\in A\inter K$.
                $x\in A$ and $x\in K$ by \obv{inter}.
                $x\in C\inter\psiCarrier{P}$ or $x\in H$
                    by \auto{union_iff}.
                \begin{byCase}
                    \caseOf{$x\in C\inter\psiCarrier{P}$.}
                        Follows by \obv{inter}.
                    \caseOf{$x\in H$.}
                        $x\in P$ by \obv{subseteq}.
                        $A\subseteq\naturals$ by
                            \obv{almost_disjoint_members_subset}.
                        Hence $x\in\naturals$ by \obv{subseteq}.
                        Hence $x\in\naturals\inter P$ by \obv{inter}.
                        $\naturals\inter P=\emptyset$ by
                            \obv{psi_carrier_naturals_disjoint_family}.
                        Contradiction by \auto{emptyset}.
                \end{byCase}
            \end{subproof}
            Hence $A\inter K\subseteq A\inter C$ by \obv{subseteq}.
            Follows by \obv{subseteq_finite_is_finite}.
        \end{subproof}
        Follows by \auto{psi_closed_iff}.
    \end{subproof}
    $K$ is closed in $\psiSpace{P}$.
    Hence $K\in\closeds{\psiSpace{P}}$ by \obv{is_closed_in}.
    $\closure{C}{\psiSpace{P}}=
        \closure{(C\inter\carrier[\psiSpace{P}])}{\psiSpace{P}}$
        by \obv{closure_trace}.
    Hence $\closure{C}{\psiSpace{P}}=
        \closure{(C\inter\psiCarrier{P})}{\psiSpace{P}}$
        by \auto{psi_space_carrier}.
    $C\inter\psiCarrier{P}\subseteq\carrier[\psiSpace{P}]$
        by \auto{psi_space_carrier}.
    Hence $\closure{(C\inter\psiCarrier{P})}{\psiSpace{P}}\subseteq K$
        by \obv{closure_is_minimal_closed_set}.
    Hence $\closure{C}{\psiSpace{P}}\subseteq K$.
    $C\inter\carrier[\psiSpace{P}]\subseteq
        \closure{C}{\psiSpace{P}}$
        by \obv{trace_is_subseteq_to_closure}.
    Hence $C\inter\psiCarrier{P}\subseteq
        \closure{C}{\psiSpace{P}}$ by \obv{psi_space_carrier}.
    Show that every element of $H$ is an element of
        $\closure{C}{\psiSpace{P}}$.
    \begin{subproof}
        Fix $A\in H$.
        $A\in P$ and $A\inter C$ is infinite.
        Suppose not.
        Hence $A\notin\closure{C}{\psiSpace{P}}$.
        $\closure{C}{\psiSpace{P}}$ is closed in $\psiSpace{P}$
            by \obv{closure_is_closed}.
        $\closure{C}{\psiSpace{P}}\subseteq\carrier[\psiSpace{P}]$
            by \obv{closed_in_subseteq_carrier}.
        Hence $\closure{C}{\psiSpace{P}}\subseteq\psiCarrier{P}$
            by \obv{psi_space_carrier}.
        Hence $A\in P\setminus\closure{C}{\psiSpace{P}}$
            by \obv{setminus}.
        Hence $A\inter\closure{C}{\psiSpace{P}}$ is finite
            by \auto{psi_closed_iff}.
        Show that every element of $A\inter C$ is an element of
            $A\inter\closure{C}{\psiSpace{P}}$.
        \begin{subproof}
            Fix $x\in A\inter C$.
            $x\in A$ and $x\in C$ by \obv{inter}.
            $A\subseteq\naturals$ by
                \obv{almost_disjoint_members_subset}.
            Hence $x\in\naturals$ by \obv{subseteq}.
            Hence $x\in\psiCarrier{P}$ by
                \auto{naturals_subset_psi_carrier,subseteq}.
            Hence $x\in C\inter\psiCarrier{P}$ by \obv{inter}.
            Hence $x\in\closure{C}{\psiSpace{P}}$ by \obv{subseteq}.
            Follows by \obv{inter}.
        \end{subproof}
        Hence $A\inter C\subseteq
            A\inter\closure{C}{\psiSpace{P}}$ by \obv{subseteq}.
        Hence $A\inter C$ is finite by \obv{subseteq_finite_is_finite}.
        Contradiction.
    \end{subproof}
    Hence $H\subseteq\closure{C}{\psiSpace{P}}$ by \obv{subseteq}.
    Hence $K\subseteq\closure{C}{\psiSpace{P}}$
        by \auto{union_iff,subseteq}.
    Follows by \auto{subseteq_antisymmetric}.
\end{proof}

\begin{corollary}\label{psi_closure_of_natural_subset}
    Suppose $P$ is an almost-disjoint family on $\naturals$.
    Suppose $B\subseteq\naturals$.
    Then
        $\closure{B}{\psiSpace{P}}=
        B\union\{A\in P\mid \text{$A\inter B$ is infinite}\}$.
\end{corollary}
\begin{proof}
    $B\subseteq\psiCarrier{P}$ by
        \auto{naturals_subset_psi_carrier,subseteq_transitive}.
    Hence $B\inter\psiCarrier{P}=B$
        by \obv{inter_absorb_supseteq_right}.
    Follows by \auto{psi_closure_formula}.
\end{proof}

\begin{corollary}\label{psi_family_point_in_closure_iff}
    Suppose $P$ is an almost-disjoint family on $\naturals$.
    Suppose $B\subseteq\naturals$ and $A\in P$.
    Then $A\in\closure{B}{\psiSpace{P}}$ iff
        $A\inter B$ is infinite.
\end{corollary}
\begin{proof}
    Show that $A\notin\naturals$.
    \begin{subproof}
        Suppose $A\in\naturals$.
        Hence $A\in\naturals\inter P$ by \obv{inter}.
        $\naturals\inter P=\emptyset$ by
            \obv{psi_carrier_naturals_disjoint_family}.
        Contradiction by \auto{emptyset}.
    \end{subproof}
    Hence $A\notin B$ by \auto{subseteq}.
    Follows by \auto{psi_closure_of_natural_subset,union_iff}.
\end{proof}

\begin{theorem}\label{psi_blocking_iff_family_closure_hit}
    Suppose $P$ is an almost-disjoint family on $\naturals$.
    Then $P$ blocks the infinite subsets of $\naturals$ iff
        for every $B\subseteq\naturals$ such that $B$ is infinite
            we have $\closure{B}{\psiSpace{P}}\inter P$ is inhabited.
\end{theorem}
\begin{proof}
    Show that if $P$ blocks the infinite subsets of $\naturals$, then
        for every $B\subseteq\naturals$ such that $B$ is infinite
            we have $\closure{B}{\psiSpace{P}}\inter P$ is inhabited.
    \begin{subproof}
        Assume $P$ blocks the infinite subsets of $\naturals$.
        Fix $B\subseteq\naturals$.
        Assume $B$ is infinite.
        Take $A\in P$ such that $A\inter B$ is infinite
            by \obv{blocks_infinite_subsets_of}.
        Hence $A\in\closure{B}{\psiSpace{P}}$ by
            \obv{psi_family_point_in_closure_iff}.
        Hence $A\in\closure{B}{\psiSpace{P}}\inter P$ by \obv{inter}.
        Follows by \auto.
    \end{subproof}
    Show that if for every $B\subseteq\naturals$ such that $B$ is infinite
            we have $\closure{B}{\psiSpace{P}}\inter P$ is inhabited,
        then $P$ blocks the infinite subsets of $\naturals$.
    \begin{subproof}
        Assume for every $B\subseteq\naturals$ such that $B$ is infinite
            we have $\closure{B}{\psiSpace{P}}\inter P$ is inhabited.
        Show that for every $B\subseteq\naturals$ such that $B$ is infinite
            there exists $A\in P$ such that $A\inter B$ is infinite.
        \begin{subproof}
            Fix $B\subseteq\naturals$.
            Assume $B$ is infinite.
            $\closure{B}{\psiSpace{P}}\inter P$ is inhabited.
            Take $A$ such that $A\in\closure{B}{\psiSpace{P}}\inter P$.
            $A\in P$ and $A\in\closure{B}{\psiSpace{P}}$ by \obv{inter}.
            Hence $A\inter B$ is infinite by
                \obv{psi_family_point_in_closure_iff}.
            Follows by \auto.
        \end{subproof}
        Follows by \obv{blocks_infinite_subsets_of}.
    \end{subproof}
    Follows by \auto.
\end{proof}

\begin{corollary}\label{psi_maximal_family_iff_every_infinite_natural_set_has_family_closure_point}
    Suppose $P$ is an almost-disjoint family on $\naturals$.
    Then $P$ is a maximal almost-disjoint family on $\naturals$ iff
        for every $B\subseteq\naturals$ such that $B$ is infinite
            we have $\closure{B}{\psiSpace{P}}\inter P$ is inhabited.
\end{corollary}
\begin{proof}
    Follows by
        \auto{maximal_almost_disjoint_iff_blocking,
            psi_blocking_iff_family_closure_hit}.
\end{proof}

\begin{proposition}\label{psi_natural_point_isolated}
    Suppose $P$ is an almost-disjoint family on $\naturals$.
    Let $n\in\naturals$.
    Then $\{n\}$ is open in $\psiSpace{P}$.
\end{proposition}
\begin{proof}
    $\{n\}\in\psiBasis{P}$ by \obv{psi_natural_singleton_in_basis}.
    $\psiBasis{P}$ is an open basis of $\psiSpace{P}$
        by \obv{psi_basis_is_open_basis}.
    Hence $\psiBasis{P}\subseteq\opens[\psiSpace{P}]$
        by \obv{open_basis}.
    Hence $\{n\}\in\opens[\psiSpace{P}]$ by \obv{subseteq}.
    Follows by \obv{is_open}.
\end{proof}

\begin{proposition}\label{psi_naturals_open}
    Suppose $P$ is an almost-disjoint family on $\naturals$.
    Then $\naturals$ is open in $\psiSpace{P}$.
\end{proposition}
\begin{proof}
    $\naturals\subseteq\psiCarrier{P}$
        by \obv{naturals_subset_psi_carrier}.
    $P\inter\naturals=\emptyset$
        by \auto{psi_carrier_naturals_disjoint_family,inter_comm}.
    For every $A\in P\inter\naturals$ we have
        $A\setminus\naturals$ is finite by \auto{emptyset}.
    Follows by \auto{psi_open_iff}.
\end{proof}

\begin{theorem}\label{psi_naturals_dense}
    Suppose $P$ is an almost-disjoint family on $\naturals$.
    Then $\naturals$ is dense in $\psiSpace{P}$.
\end{theorem}
\begin{proof}
    $\psiSpace{P}$ is a topological space
        by \obv{psi_space_is_topological_space}.
    $\naturals\subseteq\psiCarrier{P}$
        by \obv{naturals_subset_psi_carrier}.
    Hence $\naturals\subseteq\carrier[\psiSpace{P}]$
        by \auto{psi_space_carrier}.
    Hence $\naturals\inter\carrier[\psiSpace{P}]=\naturals$
        by \obv{inter_absorb_supseteq_right}.
    $\naturals\inter\carrier[\psiSpace{P}]\subseteq
        \closure{\naturals}{\psiSpace{P}}$
        by \obv{trace_is_subseteq_to_closure}.
    Hence $\naturals\subseteq\closure{\naturals}{\psiSpace{P}}$.
    Show that every element of $\psiCarrier{P}$ is an element of
        $\closure{\naturals}{\psiSpace{P}}$.
    \begin{subproof}
        Fix $x\in\psiCarrier{P}$.
        $x\in\naturals$ or $x\in P$ by \obv{psi_carrier_point_cases}.
        \begin{byCase}
            \caseOf{$x\in\naturals$.}
                Follows by \obv{subseteq}.
            \caseOf{$x\in P$.}
                $x\in P$.
                $\naturals\subseteq\naturals$ by \obv{subseteq_refl}.
                $x\subseteq\naturals$ by
                    \obv{almost_disjoint_members_subset}.
                Hence $x\inter\naturals=x$
                    by \obv{inter_absorb_supseteq_right}.
                $x$ is infinite by
                    \obv{almost_disjoint_members_infinite}.
                Hence $x\inter\naturals$ is infinite.
                Follows by \obv{psi_family_point_in_closure_iff}.
        \end{byCase}
    \end{subproof}
    Hence $\psiCarrier{P}\subseteq
        \closure{\naturals}{\psiSpace{P}}$ by \obv{subseteq}.
    $\closure{\naturals}{\psiSpace{P}}$ is closed in $\psiSpace{P}$
        by \obv{closure_is_closed}.
    Hence $\closure{\naturals}{\psiSpace{P}}\subseteq\psiCarrier{P}$
        by \auto{psi_closed_iff}.
    Hence $\closure{\naturals}{\psiSpace{P}}=\psiCarrier{P}$
        by \obv{subseteq_antisymmetric}.
    Hence $\closure{\naturals}{\psiSpace{P}}=
        \carrier[\psiSpace{P}]$ by \obv{psi_space_carrier}.
    Follows by \obv{is_dense_in}.
\end{proof}

\begin{theorem}\label{psi_family_subsets_closed}
    Suppose $P$ is an almost-disjoint family on $\naturals$.
    Suppose $Q\subseteq P$.
    Then $Q$ is closed in $\psiSpace{P}$.
\end{theorem}
\begin{proof}
    $P\subseteq\psiCarrier{P}$ by \obv{psi_family_subset_psi_carrier}.
    Hence $Q\subseteq\psiCarrier{P}$ by \obv{subseteq_transitive}.
    Show that for every $A\in P\setminus Q$ we have
        $A\inter Q$ is finite.
    \begin{subproof}
        Fix $A\in P\setminus Q$.
        $A\in P$ by \obv{setminus}.
        Show that $A\inter Q=\emptyset$.
        \begin{subproof}
            Show that every element of $A\inter Q$ is an element of
                $\emptyset$.
            \begin{subproof}
                Fix $x\in A\inter Q$.
                $x\in A$ and $x\in Q$ by \obv{inter}.
                $A\subseteq\naturals$ by
                    \obv{almost_disjoint_members_subset}.
                Hence $x\in\naturals$ by \obv{subseteq}.
                $x\in P$ by \obv{subseteq}.
                Hence $x\in\naturals\inter P$ by \obv{inter}.
                $\naturals\inter P=\emptyset$ by
                    \obv{psi_carrier_naturals_disjoint_family}.
                Contradiction by \auto{emptyset}.
            \end{subproof}
            Follows by \auto{subseteq_emptyset_iff,subseteq}.
        \end{subproof}
        Follows by \obv{emptyset_is_finite}.
    \end{subproof}
    Follows by \auto{psi_closed_iff}.
\end{proof}

\begin{proposition}\label{psi_family_closed}
    Suppose $P$ is an almost-disjoint family on $\naturals$.
    Then $P$ is closed in $\psiSpace{P}$.
\end{proposition}
\begin{proof}
    $P\subseteq P$ by \obv{subseteq_refl}.
    Follows by \obv{psi_family_subsets_closed}.
\end{proof}

\begin{theorem}\label{psi_family_discrete}
    Suppose $P$ is an almost-disjoint family on $\naturals$.
    Then $P$ is discrete in $\psiSpace{P}$.
\end{theorem}
\begin{proof}
    $\psiSpace{P}$ is a topological space
        by \obv{psi_space_is_topological_space}.
    $P\subseteq\psiCarrier{P}$ by \obv{psi_family_subset_psi_carrier}.
    Hence $P\subseteq\carrier[\psiSpace{P}]$ by \auto{psi_space_carrier}.
    Show that for every $A\in P$ there exists
        $U\in\opens[\psiSpace{P}]$ such that $U\inter P=\{A\}$.
    \begin{subproof}
        Fix $A\in P$.
        $\emptyset\subseteq\naturals$ by \obv{emptyset_subseteq}.
        $\emptyset$ is finite by \obv{emptyset_is_finite}.
        Hence $\psiNeighbourhood{A}{\emptyset}\in\psiBasis{P}$
            by \obv{psi_neighbourhood_in_basis}.
        $\psiBasis{P}$ is an open basis of $\psiSpace{P}$
            by \obv{psi_basis_is_open_basis}.
        Hence $\psiBasis{P}\subseteq\opens[\psiSpace{P}]$
            by \obv{open_basis}.
        Hence $\psiNeighbourhood{A}{\emptyset}\in
            \opens[\psiSpace{P}]$ by \obv{subseteq}.
        $\psiNeighbourhood{A}{\emptyset}\inter P=\{A\}$
            by \obv{psi_neighbourhood_family_trace}.
        Follows by \auto.
    \end{subproof}
    Follows by \auto{discrete_in_iff_ambient_open_singletons}.
\end{proof}

\begin{corollary}\label{psi_family_subset_discrete}
    Suppose $P$ is an almost-disjoint family on $\naturals$.
    Suppose $Q\subseteq P$.
    Then $Q$ is discrete in $\psiSpace{P}$.
\end{corollary}
\begin{proof}
    $P$ is discrete in $\psiSpace{P}$ by \obv{psi_family_discrete}.
    Follows by \obv{discrete_in_subset}.
\end{proof}

\begin{lemma}\label{psi_continuous_family_point_tail}
    Suppose $P$ is an almost-disjoint family on $\naturals$.
    Suppose $Y$ is a topological space.
    Suppose $f$ is continuous from $\psiSpace{P}$ to $Y$.
    Suppose $A\in P$.
    Suppose $V$ is open in $Y$ and $f(A)\in V$.
    Then there exists $F$ such that
        $F\subseteq\naturals$ and $F$ is finite and
        for every $n\in A\setminus F$ we have $f(n)\in V$.
\end{lemma}
\begin{proof}
    $\psiSpace{P}$ is a topological space
        by \obv{psi_space_is_topological_space}.
    $f\in\funs{\carrier[\psiSpace{P}]}{\carrier[Y]}$
        by \obv{continuous_map}.
    $P\subseteq\psiCarrier{P}$ by \obv{psi_family_subset_psi_carrier}.
    Hence $A\in\psiCarrier{P}$ by \obv{subseteq}.
    Hence $A\in\carrier[\psiSpace{P}]$ by \auto{psi_space_carrier}.
    $V\in\opens[Y]$ by \obv{is_open}.
    Take $U\in\opens[\psiSpace{P}]$ such that $A\in U$ and
        for every $x\in U$ we have $f(x)\in V$
        by \auto{continuous_iff_open_local}.
    $U$ is open in $\psiSpace{P}$ by \obv{is_open}.
    $A\in P\inter U$ by \obv{inter}.
    Take $F$ such that $F\subseteq\naturals$ and $F$ is finite and
        $\psiNeighbourhood{A}{F}\subseteq U$
        by \obv{psi_family_point_open_refinement}.
    Show that for every $n\in A\setminus F$ we have $f(n)\in V$.
    \begin{subproof}
        Fix $n\in A\setminus F$.
        Hence $n\in\psiNeighbourhood{A}{F}$
            by \obv{psi_neighbourhood,union_intro_right}.
        Hence $n\in U$ by \obv{subseteq}.
    \end{subproof}
    Follows by \auto.
\end{proof}

\begin{theorem}\label{psi_natural_singleton_clopen}
    Suppose $P$ is an almost-disjoint family on $\naturals$.
    Let $n\in\naturals$.
    Then $\{n\}$ is clopen in $\psiSpace{P}$.
\end{theorem}
\begin{proof}
    $\{n\}$ is open in $\psiSpace{P}$ by \obv{psi_natural_point_isolated}.
    $n\in\psiCarrier{P}$ by
        \auto{naturals_subset_psi_carrier,subseteq}.
    Hence $\{n\}\subseteq\psiCarrier{P}$
        by \obv{singleton_subset_intro}.
    Show that for every $A\in P\setminus\{n\}$ we have
        $A\inter\{n\}$ is finite.
    \begin{subproof}
        Fix $A\in P\setminus\{n\}$.
        $A\inter\{n\}\subseteq\{n\}$ by \obv{inter_lower_right}.
        $\{n\}$ is finite by \obv{singleton_is_finite}.
        Follows by \obv{subseteq_finite_is_finite}.
    \end{subproof}
    Hence $\{n\}$ is closed in $\psiSpace{P}$ by \auto{psi_closed_iff}.
    Follows by \obv{is_clopen_in}.
\end{proof}

\begin{theorem}\label{psi_neighbourhood_clopen}
    Suppose $P$ is an almost-disjoint family on $\naturals$.
    Suppose $A\in P$.
    Suppose $F\subseteq\naturals$ and $F$ is finite.
    Then $\psiNeighbourhood{A}{F}$ is clopen in $\psiSpace{P}$.
\end{theorem}
\begin{proof}
    $\psiNeighbourhood{A}{F}\in\psiBasis{P}$
        by \obv{psi_neighbourhood_in_basis}.
    $\psiBasis{P}$ is an open basis of $\psiSpace{P}$
        by \obv{psi_basis_is_open_basis}.
    Hence $\psiBasis{P}\subseteq\opens[\psiSpace{P}]$
        by \obv{open_basis}.
    Hence $\psiNeighbourhood{A}{F}\in\opens[\psiSpace{P}]$
        by \obv{subseteq}.
    Hence $\psiNeighbourhood{A}{F}$ is open in $\psiSpace{P}$
        by \obv{is_open}.
    $\psiNeighbourhood{A}{F}\subseteq\psiCarrier{P}$
        by \obv{psi_neighbourhood_subset_carrier}.
    Show that for every $B\in P\setminus\psiNeighbourhood{A}{F}$
        we have $B\inter\psiNeighbourhood{A}{F}$ is finite.
    \begin{subproof}
        Fix $B\in P\setminus\psiNeighbourhood{A}{F}$.
        $B\in P$ and $B\notin\psiNeighbourhood{A}{F}$ by \obv{setminus}.
        $A\in\psiNeighbourhood{A}{F}$
            by \obv{psi_neighbourhood_contains_center}.
        Hence $B\neq A$.
        Show that every element of $B\inter\psiNeighbourhood{A}{F}$
            is an element of $B\inter A$.
        \begin{subproof}
            Fix $x\in B\inter\psiNeighbourhood{A}{F}$.
            $x\in B$ and $x\in\psiNeighbourhood{A}{F}$ by \obv{inter}.
            $x=A$ or $x\in A\setminus F$
                by \auto{psi_neighbourhood,union_iff,singleton_iff}.
            \begin{byCase}
                \caseOf{$x=A$.}
                    $B\subseteq\naturals$ by
                        \obv{almost_disjoint_members_subset}.
                    Hence $A\in\naturals$ by \auto{subseteq}.
                    Hence $A\in\naturals\inter P$ by \obv{inter}.
                    $\naturals\inter P=\emptyset$ by
                        \obv{psi_carrier_naturals_disjoint_family}.
                    Contradiction by \auto{emptyset}.
                \caseOf{$x\in A\setminus F$.}
                    Follows by \obv{setminus,inter}.
            \end{byCase}
        \end{subproof}
        Hence $B\inter\psiNeighbourhood{A}{F}\subseteq B\inter A$
            by \obv{subseteq}.
        $B\inter A$ is finite by
            \obv{almost_disjoint_distinct_intersection_finite}.
        Follows by \obv{subseteq_finite_is_finite}.
    \end{subproof}
    Hence $\psiNeighbourhood{A}{F}$ is closed in $\psiSpace{P}$
        by \auto{psi_closed_iff}.
    Follows by \obv{is_clopen_in}.
\end{proof}

\begin{theorem}\label{psi_basic_set_clopen}
    Suppose $P$ is an almost-disjoint family on $\naturals$.
    Suppose $V\in\psiBasis{P}$.
    Then $V$ is clopen in $\psiSpace{P}$.
\end{theorem}
\begin{proof}
    (There exists $n\in\naturals$ such that $V=\{n\}$) or
        (there exist $A,F$ such that $A\in P$ and
            $F\in\finiteSubsets{\naturals}$ and
            $V=\psiNeighbourhood{A}{F}$)
        by \obv{psi_basis_member_cases}.
    \begin{byCase}
        \caseOf{There exists $n\in\naturals$ such that $V=\{n\}$.}
            Take $n\in\naturals$ such that $V=\{n\}$.
            Follows by \auto{psi_natural_singleton_clopen}.
        \caseOf{There exist $A,F$ such that $A\in P$ and
                $F\in\finiteSubsets{\naturals}$ and
                $V=\psiNeighbourhood{A}{F}$.}
            Take $A,F$ such that $A\in P$ and
                $F\in\finiteSubsets{\naturals}$ and
                $V=\psiNeighbourhood{A}{F}$.
            $F\subseteq\naturals$ and $F$ is finite
                by \auto{finite_subsets_iff}.
            Follows by \auto{psi_neighbourhood_clopen}.
    \end{byCase}
\end{proof}

\begin{theorem}\label{psi_natural_singleton_compact}
    Suppose $P$ is an almost-disjoint family on $\naturals$.
    Let $n\in\naturals$.
    Then $\{n\}$ is compact in $\psiSpace{P}$.
\end{theorem}
\begin{proof}
    $\psiSpace{P}$ is a topological space
        by \obv{psi_space_is_topological_space}.
    $n\in\psiCarrier{P}$ by
        \auto{naturals_subset_psi_carrier,subseteq}.
    Hence $\{n\}\subseteq\psiCarrier{P}$
        by \obv{singleton_subset_intro}.
    Hence $\{n\}\subseteq\carrier[\psiSpace{P}]$
        by \auto{psi_space_carrier}.
    $\subspaceTopology{\psiSpace{P}}{\{n\}}$ is a topological space
        by \obv{subspace_topology_is_topological_space}.
    $\carrier[\subspaceTopology{\psiSpace{P}}{\{n\}}]=\{n\}$
        by \obv{subspace_topology_carrier_of_subset}.
    Hence $\carrier[\subspaceTopology{\psiSpace{P}}{\{n\}}]$ is finite
        by \auto{singleton_is_finite}.
    Hence $\subspaceTopology{\psiSpace{P}}{\{n\}}$ is compact
        by \obv{finite_topological_space_is_compact}.
    Follows by \obv{compact_in}.
\end{proof}

\begin{theorem}\label{psi_neighbourhood_compact}
    Suppose $P$ is an almost-disjoint family on $\naturals$.
    Suppose $A\in P$.
    Suppose $F\subseteq\naturals$ and $F$ is finite.
    Then $\psiNeighbourhood{A}{F}$ is compact in $\psiSpace{P}$.
\end{theorem}
\begin{proof}
    Let $K=\psiNeighbourhood{A}{F}$.
    Let $Y=\subspaceTopology{\psiSpace{P}}{K}$.
    $\psiSpace{P}$ is a topological space
        by \obv{psi_space_is_topological_space}.
    $K\subseteq\psiCarrier{P}$
        by \auto{psi_neighbourhood_subset_carrier}.
    Hence $K\subseteq\carrier[\psiSpace{P}]$ by \auto{psi_space_carrier}.
    $Y$ is a topological space
        by \auto{subspace_topology_is_topological_space}.
    $\carrier[Y]=K$ by \auto{subspace_topology_carrier_of_subset}.
    $A\in K$ by \obv{psi_neighbourhood_contains_center}.
    Hence $A\in\carrier[Y]$.
    Show that for every $V\in\opens[Y]$ such that $A\in V$ we have
        $\carrier[Y]\setminus V$ is finite.
    \begin{subproof}
        Fix $V\in\opens[Y]$.
        Assume $A\in V$.
        Take $U\in\opens[\psiSpace{P}]$ such that
            $V=U\inter\carrier[Y]$ by \auto{subspace_open_iff}.
        Hence $V=U\inter K$.
        Hence $A\in U$ by \auto{inter}.
        $U$ is open in $\psiSpace{P}$ by \obv{is_open}.
        Hence $A\in P\inter U$ by \obv{inter}.
        Take $G$ such that $G\subseteq\naturals$ and $G$ is finite and
            $\psiNeighbourhood{A}{G}\subseteq U$
            by \obv{psi_family_point_open_refinement}.
        Show that every element of $K\setminus V$ is an element of $G$.
        \begin{subproof}
            Fix $x\in K\setminus V$.
            Suppose not.
            Hence $x\notin G$.
            $x\in K$ by \obv{setminus}.
            Hence $x=A$ or $x\in A\setminus F$
                by \auto{psi_neighbourhood,union_iff,singleton_iff}.
            Show that $x\in\psiNeighbourhood{A}{G}$.
            \begin{subproof}
                \begin{byCase}
                    \caseOf{$x=A$.}
                        Follows by
                            \auto{psi_neighbourhood_contains_center}.
                    \caseOf{$x\in A\setminus F$.}
                        $x\in A$ by \obv{setminus}.
                        Hence $x\in A\setminus G$ by \obv{setminus}.
                        Follows by
                            \obv{psi_neighbourhood,union_intro_right}.
                \end{byCase}
            \end{subproof}
            Hence $x\in U$ by \obv{subseteq}.
            Hence $x\in U\inter K$ by \obv{inter}.
            Hence $x\in V$.
            Contradiction by \obv{setminus}.
        \end{subproof}
        Hence $K\setminus V\subseteq G$ by \obv{subseteq}.
        Hence $K\setminus V$ is finite
            by \obv{subseteq_finite_is_finite}.
    \end{subproof}
    Hence $Y$ is compact by
        \obv{cofinite_neighbourhood_point_implies_compact}.
    Follows by \auto{compact_in}.
\end{proof}

\begin{theorem}\label{psi_basic_set_compact}
    Suppose $P$ is an almost-disjoint family on $\naturals$.
    Suppose $V\in\psiBasis{P}$.
    Then $V$ is compact in $\psiSpace{P}$.
\end{theorem}
\begin{proof}
    (There exists $n\in\naturals$ such that $V=\{n\}$) or
        (there exist $A,F$ such that $A\in P$ and
            $F\in\finiteSubsets{\naturals}$ and
            $V=\psiNeighbourhood{A}{F}$)
        by \obv{psi_basis_member_cases}.
    \begin{byCase}
        \caseOf{There exists $n\in\naturals$ such that $V=\{n\}$.}
            Take $n\in\naturals$ such that $V=\{n\}$.
            Follows by \auto{psi_natural_singleton_compact}.
        \caseOf{There exist $A,F$ such that $A\in P$ and
                $F\in\finiteSubsets{\naturals}$ and
                $V=\psiNeighbourhood{A}{F}$.}
            Take $A,F$ such that $A\in P$ and
                $F\in\finiteSubsets{\naturals}$ and
                $V=\psiNeighbourhood{A}{F}$.
            $F\subseteq\naturals$ and $F$ is finite
                by \auto{finite_subsets_iff}.
            Follows by \auto{psi_neighbourhood_compact}.
    \end{byCase}
\end{proof}

\begin{theorem}\label{psi_compact_clopen_refinement}
    Suppose $P$ is an almost-disjoint family on $\naturals$.
    Suppose $U$ is open in $\psiSpace{P}$ and $x\in U$.
    Then there exists $V$ such that
        $V$ is compact in $\psiSpace{P}$ and
        $V$ is clopen in $\psiSpace{P}$ and
        $x\in V\subseteq U$.
\end{theorem}
\begin{proof}
    $U\in\opens[\psiSpace{P}]$ by \obv{is_open}.
    $\psiBasis{P}$ is an open basis of $\psiSpace{P}$
        by \obv{psi_basis_is_open_basis}.
    Take $V\in\psiBasis{P}$ such that $x\in V\subseteq U$
        by \obv{open_basis}.
    $V$ is compact in $\psiSpace{P}$ by \obv{psi_basic_set_compact}.
    $V$ is clopen in $\psiSpace{P}$ by \obv{psi_basic_set_clopen}.
    Follows by \auto.
\end{proof}

\begin{proposition}\label{psi_natural_family_neighbourhoods_disjoint}
    Suppose $P$ is an almost-disjoint family on $\naturals$.
    Suppose $n\in\naturals$ and $A\in P$.
    Then $\{n\}$ is disjoint from $\psiNeighbourhood{A}{\{n\}}$.
\end{proposition}
\begin{proof}
    Show that every element of
        $\{n\}\inter\psiNeighbourhood{A}{\{n\}}$ is an element of
        $\emptyset$.
    \begin{subproof}
        Fix $x\in\{n\}\inter\psiNeighbourhood{A}{\{n\}}$.
        $x=n$ and $x\in\psiNeighbourhood{A}{\{n\}}$
            by \auto{inter,singleton_iff}.
        Hence $x=A$ or $x\in A\setminus\{n\}$
            by \auto{psi_neighbourhood,union_iff,singleton_iff}.
        \begin{byCase}
            \caseOf{$x=A$.}
                Hence $A\in\naturals\inter P$ by \auto{inter}.
                $\naturals\inter P=\emptyset$
                    by \obv{psi_carrier_naturals_disjoint_family}.
                Contradiction by \auto{emptyset}.
            \caseOf{$x\in A\setminus\{n\}$.}
                Hence $x\notin\{n\}$ by \obv{setminus}.
                Contradiction by \obv{singleton_intro}.
        \end{byCase}
    \end{subproof}
    Hence $\{n\}\inter\psiNeighbourhood{A}{\{n\}}\subseteq\emptyset$
        by \obv{subseteq}.
    $\emptyset\subseteq
        \{n\}\inter\psiNeighbourhood{A}{\{n\}}$
        by \obv{emptyset_subseteq}.
    Hence $\{n\}\inter\psiNeighbourhood{A}{\{n\}}=\emptyset$
        by \obv{subseteq_antisymmetric}.
    Follows by \auto{disjoint,inter,emptyset}.
\end{proof}

\begin{proposition}\label{psi_distinct_family_neighbourhoods_disjoint}
    Suppose $P$ is an almost-disjoint family on $\naturals$.
    Suppose $A,B\in P$ and $A\neq B$.
    Let $H=A\inter B$.
    Then $\psiNeighbourhood{A}{H}$ is disjoint from
        $\psiNeighbourhood{B}{H}$.
\end{proposition}
\begin{proof}
    $A\subseteq\naturals$ by
        \obv{almost_disjoint_members_subset}.
    Hence $H\subseteq\naturals$ by \auto{inter,subseteq}.
    $\psiNeighbourhood{A}{H}\inter\psiNeighbourhood{B}{H}
        \subseteq H$
        by \auto{psi_neighbourhood_distinct_centers_intersection}.
    Show that every element of
        $\psiNeighbourhood{A}{H}\inter\psiNeighbourhood{B}{H}$
        is an element of $\emptyset$.
    \begin{subproof}
        Fix $x\in
            \psiNeighbourhood{A}{H}\inter\psiNeighbourhood{B}{H}$.
        Hence $x\in H$ by \obv{subseteq}.
        Hence $x\in\naturals$ by \obv{subseteq}.
        $x\in\psiNeighbourhood{A}{H}$ by \obv{inter}.
        Hence $x=A$ or $x\in A\setminus H$
            by \auto{psi_neighbourhood,union_iff,singleton_iff}.
        Show that $x\neq A$.
        \begin{subproof}
            Suppose $x=A$.
            Hence $A\in\naturals\inter P$ by \auto{inter}.
            $\naturals\inter P=\emptyset$
                by \obv{psi_carrier_naturals_disjoint_family}.
            Contradiction by \auto{emptyset}.
        \end{subproof}
        Hence $x\in A\setminus H$.
        Hence $x\notin H$ by \obv{setminus}.
        Contradiction.
    \end{subproof}
    Hence $\psiNeighbourhood{A}{H}\inter
        \psiNeighbourhood{B}{H}\subseteq\emptyset$ by \obv{subseteq}.
    $\emptyset\subseteq\psiNeighbourhood{A}{H}\inter
        \psiNeighbourhood{B}{H}$ by \obv{emptyset_subseteq}.
    Hence $\psiNeighbourhood{A}{H}\inter
        \psiNeighbourhood{B}{H}=\emptyset$
        by \obv{subseteq_antisymmetric}.
    Follows by \auto{disjoint,inter,emptyset}.
\end{proof}

\begin{theorem}\label{psi_space_is_hausdorff}
    Suppose $P$ is an almost-disjoint family on $\naturals$.
    Then $\psiSpace{P}$ is Hausdorff.
\end{theorem}
\begin{proof}
    $\psiSpace{P}$ is a topological space
        by \obv{psi_space_is_topological_space}.
    Show that for all $x,y\in\carrier[\psiSpace{P}]$ such that
        $x\neq y$ there exist $U,V\in\opens[\psiSpace{P}]$ such that
        $x\in U$ and $y\in V$ and $U$ is disjoint from $V$.
    \begin{subproof}
        Fix $x,y\in\carrier[\psiSpace{P}]$.
        Assume $x\neq y$.
        $x\in\psiCarrier{P}$ and $y\in\psiCarrier{P}$
            by \auto{psi_space_carrier}.
        Hence ($x\in\naturals$ and $y\in\naturals$) or
            ($x\in\naturals$ and $y\in P$) or
            ($x\in P$ and $y\in\naturals$) or
            ($x\in P$ and $y\in P$)
            by \obv{psi_carrier_point_cases}.
        Show that there exist $U,V\in\opens[\psiSpace{P}]$ such that
            $x\in U$ and $y\in V$ and $U$ is disjoint from $V$.
        \begin{subproof}
            \begin{byCase}
                \caseOf{$x\in\naturals$ and $y\in\naturals$.}
                    $\{x\}$ is open in $\psiSpace{P}$ and
                        $\{y\}$ is open in $\psiSpace{P}$
                        by \obv{psi_natural_point_isolated}.
                    Hence $\{x\},\{y\}\in\opens[\psiSpace{P}]$
                        by \obv{is_open}.
                    $x\in\{x\}$ and $y\in\{y\}$
                        by \obv{singleton_intro}.
                    $\{x\}\inter\{y\}=\emptyset$
                        by \auto{setext,inter,singleton_iff,emptyset}.
                    Hence $\{x\}$ is disjoint from $\{y\}$
                        by \auto{disjoint,inter,emptyset}.
                    Follows by \auto.
                \caseOf{$x\in\naturals$ and $y\in P$.}
                    $\{x\}$ is open in $\psiSpace{P}$
                        by \obv{psi_natural_point_isolated}.
                    $\psiNeighbourhood{y}{\{x\}}$ is open in
                        $\psiSpace{P}$ by
                        \auto{psi_neighbourhood_clopen,singleton_subset_intro,
                            singleton_is_finite,is_clopen_in}.
                    Hence $\{x\},\psiNeighbourhood{y}{\{x\}}
                        \in\opens[\psiSpace{P}]$ by \obv{is_open}.
                    $x\in\{x\}$ by \obv{singleton_intro}.
                    $y\in\psiNeighbourhood{y}{\{x\}}$ by
                        \obv{psi_neighbourhood_contains_center}.
                    $\{x\}$ is disjoint from
                        $\psiNeighbourhood{y}{\{x\}}$ by
                        \obv{psi_natural_family_neighbourhoods_disjoint}.
                    Follows by \auto.
                \caseOf{$x\in P$ and $y\in\naturals$.}
                    $\{y\}$ is open in $\psiSpace{P}$
                        by \obv{psi_natural_point_isolated}.
                    $\psiNeighbourhood{x}{\{y\}}$ is open in
                        $\psiSpace{P}$ by
                        \auto{psi_neighbourhood_clopen,singleton_subset_intro,
                            singleton_is_finite,is_clopen_in}.
                    Hence $\psiNeighbourhood{x}{\{y\}},\{y\}
                        \in\opens[\psiSpace{P}]$ by \obv{is_open}.
                    $x\in\psiNeighbourhood{x}{\{y\}}$ by
                        \obv{psi_neighbourhood_contains_center}.
                    $y\in\{y\}$ by \obv{singleton_intro}.
                    $\{y\}$ is disjoint from
                        $\psiNeighbourhood{x}{\{y\}}$ by
                        \obv{psi_natural_family_neighbourhoods_disjoint}.
                    Hence $\psiNeighbourhood{x}{\{y\}}$ is disjoint
                        from $\{y\}$ by \obv{disjoint_symmetric}.
                    Follows by \auto.
                \caseOf{$x\in P$ and $y\in P$.}
                    Let $H=x\inter y$.
                    $H$ is finite by
                        \auto{almost_disjoint_distinct_intersection_finite}.
                    $x\subseteq\naturals$ by
                        \obv{almost_disjoint_members_subset}.
                    Hence $H\subseteq\naturals$ by \auto{inter,subseteq}.
                    $\psiNeighbourhood{x}{H}$ and
                        $\psiNeighbourhood{y}{H}$ are open in
                        $\psiSpace{P}$ by \auto{psi_neighbourhood_clopen,
                            is_clopen_in}.
                    Hence $\psiNeighbourhood{x}{H},
                        \psiNeighbourhood{y}{H}\in
                        \opens[\psiSpace{P}]$ by \obv{is_open}.
                    $x\in\psiNeighbourhood{x}{H}$ and
                        $y\in\psiNeighbourhood{y}{H}$ by
                        \obv{psi_neighbourhood_contains_center}.
                    $\psiNeighbourhood{x}{H}$ is disjoint from
                        $\psiNeighbourhood{y}{H}$ by
                        \obv{psi_distinct_family_neighbourhoods_disjoint}.
                    Follows by \auto.
            \end{byCase}
        \end{subproof}
    \end{subproof}
    Follows by \obv{is_hausdorff}.
\end{proof}

\begin{corollary}\label{psi_space_is_teeone}
    Suppose $P$ is an almost-disjoint family on $\naturals$.
    Then $\psiSpace{P}$ is \teeone.
\end{corollary}
\begin{proof}
    $\psiSpace{P}$ is Hausdorff by \obv{psi_space_is_hausdorff}.
    Show that for all $x,y\in\carrier[\psiSpace{P}]$ such that
        $x\neq y$ there exists $U\in\opens[\psiSpace{P}]$ such that
        $x\in U\not\ni y$.
    \begin{subproof}
        Fix $x,y\in\carrier[\psiSpace{P}]$.
        Assume $x\neq y$.
        Take $U,V\in\opens[\psiSpace{P}]$ such that
            $x\in U$ and $y\in V$ and $U$ is disjoint from $V$
            by \obv{is_hausdorff}.
        Show that $y\notin U$.
        \begin{subproof}
            Suppose $y\in U$.
            Contradiction by \auto{disjoint}.
        \end{subproof}
        Follows by \auto.
    \end{subproof}
    Follows by \obv{teeone}.
\end{proof}

\begin{corollary}\label{psi_singletons_closed}
    Suppose $P$ is an almost-disjoint family on $\naturals$.
    Let $x\in\psiCarrier{P}$.
    Then $\{x\}$ is closed in $\psiSpace{P}$.
\end{corollary}
\begin{proof}
    $\psiSpace{P}$ is a topological space
        by \obv{psi_space_is_topological_space}.
    $x\in\carrier[\psiSpace{P}]$ by \auto{psi_space_carrier}.
    $\psiSpace{P}$ is \teeone by \obv{psi_space_is_teeone}.
    Follows by \auto{teeone_iff_singletons_closed}.
\end{proof}

\begin{theorem}\label{psi_space_has_clopen_basis}
    Suppose $P$ is an almost-disjoint family on $\naturals$.
    Then $\psiSpace{P}$ has a clopen basis.
\end{theorem}
\begin{proof}
    $\psiBasis{P}$ is an open basis of $\psiSpace{P}$
        by \obv{psi_basis_is_open_basis}.
    For every $V\in\psiBasis{P}$ we have $V$ is clopen in
        $\psiSpace{P}$ by \obv{psi_basic_set_clopen}.
    Follows by \obv{clopen_open_basis_implies_has_clopen_basis}.
\end{proof}

\begin{theorem}\label{psi_space_is_completely_regular}
    Suppose $P$ is an almost-disjoint family on $\naturals$.
    Then $\psiSpace{P}$ is completely regular.
\end{theorem}
\begin{proof}
    $\psiSpace{P}$ has a clopen basis
        by \obv{psi_space_has_clopen_basis}.
    Follows by \obv{clopen_basis_implies_completely_regular}.
\end{proof}

\begin{theorem}\label{psi_space_is_tychonoff}
    Suppose $P$ is an almost-disjoint family on $\naturals$.
    Then $\psiSpace{P}$ is Tychonoff.
\end{theorem}
\begin{proof}
    $\psiSpace{P}$ is completely regular
        by \obv{psi_space_is_completely_regular}.
    $\psiSpace{P}$ is Hausdorff by \obv{psi_space_is_hausdorff}.
    Follows by \obv{is_tychonoff}.
\end{proof}

\begin{theorem}\label{psi_space_is_locally_compact}
    Suppose $P$ is an almost-disjoint family on $\naturals$.
    Then $\psiSpace{P}$ is locally compact.
\end{theorem}
\begin{proof}
    $\psiBasis{P}$ is an open basis of $\psiSpace{P}$
        by \obv{psi_basis_is_open_basis}.
    For every $V\in\psiBasis{P}$ we have
        $V$ is compact in $\psiSpace{P}$ by \obv{psi_basic_set_compact}.
    Follows by \obv{compact_open_basis_implies_locally_compact}.
\end{proof}

\begin{theorem}\label{psi_space_is_separable}
    Suppose $P$ is an almost-disjoint family on $\naturals$.
    Then $\psiSpace{P}$ is separable.
\end{theorem}
\begin{proof}
    $\psiSpace{P}$ is a topological space
        by \obv{psi_space_is_topological_space}.
    $\naturals$ is countable by \obv{naturals_is_countable}.
    $\naturals$ is dense in $\psiSpace{P}$ by \obv{psi_naturals_dense}.
    Follows by \auto{is_separable}.
\end{proof}

\begin{definition}\label{psi_standard_local_basis_map}
    $\psiStandardLocalBasisMap{A}=
        \{(n,\psiNeighbourhood{A}{n})\mid n\in\naturals\}$.
\end{definition}

\begin{definition}\label{psi_standard_local_basis}
    $\psiStandardLocalBasis{A}=\ran{\psiStandardLocalBasisMap{A}}$.
\end{definition}

\begin{proposition}\label{psi_standard_local_basis_map_membership}
    $(n,V)\in\psiStandardLocalBasisMap{A}$ iff
        $n\in\naturals$ and $V=\psiNeighbourhood{A}{n}$.
\end{proposition}
\begin{proof}
    Follows by \auto{psi_standard_local_basis_map,pair_eq_iff}.
\end{proof}

\begin{proposition}\label{psi_standard_local_basis_map_type}
    $\psiStandardLocalBasisMap{A}$ is a function on $\naturals$.
\end{proposition}
\begin{proof}
    Let $b=\psiStandardLocalBasisMap{A}$.
    $b$ is a relation by
        \auto{relation,psi_standard_local_basis_map}.
    $b$ is right-unique by
        \auto{rightunique,psi_standard_local_basis_map_membership}.
    $\dom{b}=\naturals$ by
        \auto{setext,dom_iff,psi_standard_local_basis_map_membership}.
\end{proof}

\begin{proposition}\label{psi_standard_local_basis_map_application}
    Let $n\in\naturals$.
    Then $(\psiStandardLocalBasisMap{A})(n)=
        \psiNeighbourhood{A}{n}$.
\end{proposition}
\begin{proof}
    $\psiStandardLocalBasisMap{A}$ is a function on $\naturals$
        by \obv{psi_standard_local_basis_map_type}.
    $(n,\psiNeighbourhood{A}{n})\in\psiStandardLocalBasisMap{A}$
        by \auto{psi_standard_local_basis_map_membership}.
    Follows by \obv{function_apply_intro}.
\end{proof}

\begin{proposition}\label{psi_standard_local_basis_members}
    $V\in\psiStandardLocalBasis{A}$ iff there exists $n\in\naturals$
        such that $V=\psiNeighbourhood{A}{n}$.
\end{proposition}
\begin{proof}
    $\psiStandardLocalBasisMap{A}$ is a function on $\naturals$
        by \obv{psi_standard_local_basis_map_type}.
    Show that if $V\in\psiStandardLocalBasis{A}$, then there exists
        $n\in\naturals$ such that $V=\psiNeighbourhood{A}{n}$.
    \begin{subproof}
        Assume $V\in\psiStandardLocalBasis{A}$.
        Hence $V\in\ran{\psiStandardLocalBasisMap{A}}$
            by \obv{psi_standard_local_basis}.
        Take $n\in\dom{\psiStandardLocalBasisMap{A}}$ such that
            $(\psiStandardLocalBasisMap{A})(n)=V$
            by \auto{fun_ran_iff}.
        Hence $n\in\naturals$.
        $(\psiStandardLocalBasisMap{A})(n)=
            \psiNeighbourhood{A}{n}$ by
            \obv{psi_standard_local_basis_map_application}.
        Follows by \auto.
    \end{subproof}
    Show that if there exists $n\in\naturals$ such that
        $V=\psiNeighbourhood{A}{n}$, then
        $V\in\psiStandardLocalBasis{A}$.
    \begin{subproof}
        Assume there exists $n\in\naturals$ such that
            $V=\psiNeighbourhood{A}{n}$.
        Take $n\in\naturals$ such that
            $V=\psiNeighbourhood{A}{n}$.
        $n\in\dom{\psiStandardLocalBasisMap{A}}$.
        $(\psiStandardLocalBasisMap{A})(n)=V$ by
            \auto{psi_standard_local_basis_map_application}.
        Hence $V\in\ran{\psiStandardLocalBasisMap{A}}$
            by \auto{fun_ran_iff}.
        Follows by \obv{psi_standard_local_basis}.
    \end{subproof}
    Follows by \auto.
\end{proof}

\begin{theorem}\label{psi_standard_local_basis_is_countable}
    $\psiStandardLocalBasis{A}$ is countable.
\end{theorem}
\begin{proof}
    $\psiStandardLocalBasisMap{A}$ is a function on $\naturals$
        by \obv{psi_standard_local_basis_map_type}.
    Hence $\psiStandardLocalBasisMap{A}$ is a function and
        $\dom{\psiStandardLocalBasisMap{A}}=\naturals$.
    $\naturals$ is countable by \obv{naturals_is_countable}.
    $\naturals\subseteq\dom{\psiStandardLocalBasisMap{A}}$
        by \obv{subseteq_refl}.
    Hence $\img{\psiStandardLocalBasisMap{A}}{\naturals}$ is countable
        by \obv{countable_function_image}.
    $\img{\psiStandardLocalBasisMap{A}}{\naturals}=
        \ran{\psiStandardLocalBasisMap{A}}$ by \auto{img_dom}.
    Follows by \auto{psi_standard_local_basis}.
\end{proof}

\begin{theorem}\label{psi_standard_local_basis_is_local_basis}
    Suppose $P$ is an almost-disjoint family on $\naturals$.
    Suppose $A\in P$.
    Suppose $A\in\carrier[\psiSpace{P}]$.
    Then $\psiStandardLocalBasis{A}$ is a local basis at $A$ in
        $\psiSpace{P}$.
\end{theorem}
\begin{proof}
    $\psiSpace{P}$ is a topological space
        by \obv{psi_space_is_topological_space}.
    $A\in\psiCarrier{P}$ by
        \auto{psi_family_subset_psi_carrier,subseteq}.
    Hence $A\in\carrier[\psiSpace{P}]$ by \obv{psi_space_carrier}.
    Show that for every $V\in\psiStandardLocalBasis{A}$ we have
        $V\in\opens[\psiSpace{P}]$.
    \begin{subproof}
        Fix $V\in\psiStandardLocalBasis{A}$.
        Take $n\in\naturals$ such that
            $V=\psiNeighbourhood{A}{n}$
            by \auto{psi_standard_local_basis_members}.
        $n$ is finite by \obv{natural_number_is_finite}.
        $n\subseteq\naturals$ by
            \obv{naturals_is_transitive_set,transitiveset_iff_subseteq}.
        Hence $V\in\psiBasis{P}$ by
            \auto{psi_neighbourhood_in_basis}.
        $\psiBasis{P}$ is an open basis of $\psiSpace{P}$
            by \obv{psi_basis_is_open_basis}.
        Hence $\psiBasis{P}\subseteq\opens[\psiSpace{P}]$
            by \obv{open_basis}.
        Follows by \obv{subseteq}.
    \end{subproof}
    Hence $\psiStandardLocalBasis{A}\subseteq
        \opens[\psiSpace{P}]$ by \obv{subseteq}.
    For every $V\in\psiStandardLocalBasis{A}$ we have $A\in V$
        by \auto{psi_standard_local_basis_members,
            psi_neighbourhood_contains_center}.
    Show that for all $U$ such that
        $U\in\opens[\psiSpace{P}]$ and $A\in U$
        there exists $V\in\psiStandardLocalBasis{A}$ such that
            $V\subseteq U$.
    \begin{subproof}
        Fix $U$.
        Assume $U\in\opens[\psiSpace{P}]$ and $A\in U$.
        $U$ is open in $\psiSpace{P}$ by \obv{is_open}.
        Hence $A\in P\inter U$ by \obv{inter}.
        Take $F$ such that $F\subseteq\naturals$ and $F$ is finite and
            $\psiNeighbourhood{A}{F}\subseteq U$
            by \obv{psi_family_point_open_refinement}.
        Take $n\in\naturals$ such that $F\subseteq n$
            by \obv{finite_subset_of_naturals_bounded}.
        $\psiNeighbourhood{A}{n}\subseteq\psiNeighbourhood{A}{F}$
            by \obv{psi_neighbourhood_antitone}.
        Hence $\psiNeighbourhood{A}{n}\subseteq U$
            by \obv{subseteq_transitive}.
        $\psiNeighbourhood{A}{n}\in\psiStandardLocalBasis{A}$
            by \auto{psi_standard_local_basis_members}.
        Follows by \auto.
    \end{subproof}
    Follows by \obv{local_basis}.
\end{proof}

\begin{theorem}\label{psi_space_is_first_countable}
    Suppose $P$ is an almost-disjoint family on $\naturals$.
    Then $\psiSpace{P}$ is first-countable.
\end{theorem}
\begin{proof}
    $\psiSpace{P}$ is a topological space
        by \obv{psi_space_is_topological_space}.
    Show that for every $x\in\carrier[\psiSpace{P}]$ there exists
        a countable set $B$ such that
            $B$ is a local basis at $x$ in $\psiSpace{P}$.
    \begin{subproof}
        Fix $x\in\carrier[\psiSpace{P}]$.
        Hence $x\in\psiCarrier{P}$ by \auto{psi_space_carrier}.
        Hence $x\in\naturals$ or $x\in P$
            by \obv{psi_carrier_point_cases}.
        Show that there exists a countable set $B$ such that
            $B$ is a local basis at $x$ in $\psiSpace{P}$.
        \begin{subproof}
            \begin{byCase}
                \caseOf{$x\in\naturals$.}
                    Let $B=\{\{x\}\}$.
                    $B$ is finite by \auto{singleton_is_finite}.
                    Hence $B$ is countable by \obv{finite_is_countable}.
                    $\{x\}$ is open in $\psiSpace{P}$ by
                        \obv{psi_natural_point_isolated}.
                    Hence $\{x\}\in\opens[\psiSpace{P}]$
                        by \obv{is_open}.
                    Hence $B\subseteq\opens[\psiSpace{P}]$
                        by \auto{singleton_subset_intro}.
                    For every $V\in B$ we have $x\in V$
                        by \auto{singleton_iff,singleton_intro}.
                    Show that for all $U$ such that
                        $U\in\opens[\psiSpace{P}]$ and $x\in U$
                        there exists $V\in B$ such that $V\subseteq U$.
                    \begin{subproof}
                        Fix $U$.
                        Assume $U\in\opens[\psiSpace{P}]$ and $x\in U$.
                        $\{x\}\in B$ by \obv{singleton_intro}.
                        $\{x\}\subseteq U$ by
                            \obv{singleton_subset_intro}.
                        Follows by \auto.
                    \end{subproof}
                    Hence $B$ is a local basis at $x$ in
                        $\psiSpace{P}$ by \obv{local_basis}.
                    Follows by \auto.
                \caseOf{$x\in P$.}
                    $\psiStandardLocalBasis{x}$ is countable by
                        \obv{psi_standard_local_basis_is_countable}.
                    $\psiStandardLocalBasis{x}$ is a local basis at $x$
                        in $\psiSpace{P}$ by
                        \obv{psi_standard_local_basis_is_local_basis}.
                    Follows by \auto.
            \end{byCase}
        \end{subproof}
    \end{subproof}
    Follows by \obv{is_first_countable}.
\end{proof}

\begin{theorem}\label{infinite_psi_space_is_not_countably_compact}
    Suppose $P$ is an almost-disjoint family on $\naturals$.
    Suppose $P$ is infinite.
    Then $\psiSpace{P}$ is not countably compact.
\end{theorem}
\begin{proof}
    Take $C\subseteq P$ such that $C$ is countably infinite
        by \obv{infinite_set_has_countably_infinite_subset}.
    $C$ is closed in $\psiSpace{P}$ by \obv{psi_family_subsets_closed}.
    $C$ is discrete in $\psiSpace{P}$
        by \obv{psi_family_subset_discrete}.
    Follows by
        \obv{countably_infinite_closed_discrete_obstructs_countable_compactness}.
\end{proof}

\begin{corollary}\label{infinite_almost_disjoint_psi_space_is_not_countably_compact}
    Suppose $P$ is an infinite almost-disjoint family on $\naturals$.
    Then $\psiSpace{P}$ is not countably compact.
\end{corollary}
\begin{proof}
    $P$ is an almost-disjoint family on $\naturals$.
    $P$ is infinite.
    Follows by \obv{infinite_psi_space_is_not_countably_compact}.
\end{proof}

\begin{corollary}\label{infinite_almost_disjoint_psi_space_is_not_compact}
    Suppose $P$ is an infinite almost-disjoint family on $\naturals$.
    Then $\psiSpace{P}$ is not compact.
\end{corollary}
\begin{proof}
    $P$ is an almost-disjoint family on $\naturals$.
    $\psiSpace{P}$ is not countably compact by
        \obv{infinite_almost_disjoint_psi_space_is_not_countably_compact}.
    Suppose $\psiSpace{P}$ is compact.
    Hence $\psiSpace{P}$ is countably compact by
        \obv{compact_implies_countably_compact}.
    Contradiction.
\end{proof}

\begin{lemma}\label{psi_continuous_real_map_eventually_unit_close}
    Suppose $P$ is an almost-disjoint family on $\naturals$.
    Suppose $f$ is continuous from $\psiSpace{P}$ to $\RealTopology$.
    Suppose $A\in P$.
    Then there exists $F$ such that $F\subseteq\naturals$ and
        $F$ is finite and for every $n\in A\setminus F$ we have
            $\surrealLt{\metricDistance[\RealMetric](f(A),f(n))}
                {\surrealOne}$.
\end{lemma}
\begin{proof}
    $f\in\funs{\carrier[\psiSpace{P}]}{\carrier[\RealTopology]}$
        by \obv{continuous_map}.
    $A\in\carrier[\psiSpace{P}]$ by
        \auto{psi_family_subset_psi_carrier,psi_space_carrier,subseteq}.
    Hence $f(A)\in\carrier[\RealTopology]$ by \obv{funs_type_apply}.
    Hence $f(A)\in\reals$ by \obv{real_topology_carrier}.
    $\metricBall{\RealMetric}{f(A)}{\surrealOne}
        \in\opens[\RealTopology]$ and
        $f(A)\in\metricBall{\RealMetric}{f(A)}{\surrealOne}$ by
        \obv{real_unit_ball_is_open_neighbourhood}.
    Hence $\metricBall{\RealMetric}{f(A)}{\surrealOne}$ is open in
        $\RealTopology$ by \obv{is_open}.
    Take $F$ such that $F\subseteq\naturals$ and $F$ is finite and
        for every $n\in A\setminus F$ we have
            $f(n)\in\metricBall{\RealMetric}{f(A)}{\surrealOne}$
        by \obv{psi_continuous_family_point_tail}.
    Show that for every $n\in A\setminus F$ we have
        $\surrealLt{\metricDistance[\RealMetric](f(A),f(n))}
            {\surrealOne}$.
    \begin{subproof}
        Fix $n\in A\setminus F$.
        $f(n)\in\metricBall{\RealMetric}{f(A)}{\surrealOne}$.
        Follows by \obv{metric_ball}.
    \end{subproof}
    Follows by \auto.
\end{proof}

\begin{lemma}\label{psi_continuous_real_map_bounded_if_naturals_bounded}
    Suppose $P$ is an almost-disjoint family on $\naturals$.
    Suppose $f$ is continuous from $\psiSpace{P}$ to $\RealTopology$.
    Suppose $\img{f}{\naturals}$ is real bounded.
    Then $\ran{f}$ is real bounded.
\end{lemma}
\begin{proof}
    $f\in\funs{\carrier[\psiSpace{P}]}{\carrier[\RealTopology]}$
        by \obv{continuous_map}.
    Hence $f\in\funs{\psiCarrier{P}}{\reals}$ by
        \auto{psi_space_carrier,real_topology_carrier}.
    Hence $f$ is a function and $\dom{f}=\psiCarrier{P}$ by
        \obv{funs_is_function,funs}.
    Take $M\in\reals$ such that $\surrealLe{\surrealZero}{M}$ and
        for every $y\in\img{f}{\naturals}$ we have
            $\surrealLe{\abs{y}}{M}$ by \obv{real_bounded}.
    Let $K=\surrealAdd{M}{\surrealOne}$.
    $\surrealOne\in\reals$ by \obv{surreal_one_is_real}.
    Hence $K\in\reals$ by \auto{surreal_real_addition_closed}.
    $\surrealLe{\surrealZero}{\surrealOne}$ by
        \obv{surreal_zero_lt_one,surreal_le}.
    Hence $\surrealLe{
        \surrealAdd{\surrealZero}{\surrealZero}}{K}$ by
        \auto{surreal_add_le_both,surreal_real_is_surreal,
            surreal_one_is_surreal,empty_surreal}.
    $\surrealAdd{\surrealZero}{\surrealZero}=\surrealZero$ by
        \obv{surreal_add_zero_left,empty_surreal}.
    Hence $\surrealLe{\surrealZero}{K}$ by \auto.
    $\ran{f}\subseteq\reals$ by \obv{funs_ran}.
    Show that for every $y\in\ran{f}$ we have
        $\surrealLe{\abs{y}}{K}$.
    \begin{subproof}
        Fix $y\in\ran{f}$.
        Take $x\in\dom{f}$ such that $f(x)=y$ by \auto{fun_ran_iff}.
        Hence $x\in\psiCarrier{P}$.
        Hence $x\in\naturals$ or $x\in P$
            by \obv{psi_carrier_point_cases}.
        Show that $\surrealLe{\abs{y}}{K}$.
        \begin{subproof}
            \begin{byCase}
                \caseOf{$x\in\naturals$.}
                    $\naturals\subseteq\psiCarrier{P}$ by
                        \obv{naturals_subset_psi_carrier}.
                    Hence $x\in\dom{f}\inter\naturals$ by
                        \obv{inter,subseteq}.
                    Hence $f(x)\in\img{f}{\naturals}$ by
                        \obv{img_of_function_intro}.
                    Hence $\surrealLe{\abs{f(x)}}{M}$.
                    $\surrealLe{\surrealAdd{M}{\surrealZero}}{K}$ by
                        \auto{surreal_add_le_right,surreal_real_is_surreal,
                            surreal_one_is_surreal,empty_surreal}.
                    $\surrealAdd{M}{\surrealZero}=M$ by
                        \obv{surreal_add_zero_right,
                            surreal_real_is_surreal}.
                    Hence $\surrealLe{M}{K}$ by \auto.
                    $f(x)\in\reals$ by \obv{funs_type_apply}.
                    Hence $\surreal{\abs{f(x)}}$ and $\surreal{M}$ and
                        $\surreal{K}$ by
                        \obv{surreal_real_is_surreal,
                            surreal_abs_is_surreal}.
                    Hence $\surrealLe{\abs{f(x)}}{K}$ by
                        \auto{surreal_le_transitive}.
                    Follows by \auto.
                \caseOf{$x\in P$.}
                    Take $F$ such that $F\subseteq\naturals$ and
                        $F$ is finite and for every $n\in x\setminus F$
                        we have
                            $\surrealLt{
                                \metricDistance[\RealMetric](f(x),f(n))}
                                {\surrealOne}$ by
                        \obv{psi_continuous_real_map_eventually_unit_close}.
                    $x$ is infinite by
                        \obv{almost_disjoint_members_infinite}.
                    $x\setminus F$ is inhabited by
                        \obv{infinite_minus_finite_inhabited}.
                    Take $n$ such that $n\in x\setminus F$.
                    Hence $n\in x$ by \obv{setminus}.
                    $x\subseteq\naturals$ by
                        \obv{almost_disjoint_members_subset}.
                    Hence $n\in\naturals$ by \obv{subseteq}.
                    $\naturals\subseteq\psiCarrier{P}$ by
                        \obv{naturals_subset_psi_carrier}.
                    Hence $n\in\dom{f}\inter\naturals$ by
                        \auto{inter,subseteq}.
                    Hence $f(n)\in\img{f}{\naturals}$ by
                        \obv{img_of_function_intro}.
                    Hence $\surrealLe{\abs{f(n)}}{M}$.
                    $f(x)\in\reals$ by \obv{funs_type_apply}.
                    $n\in\psiCarrier{P}$ by \obv{subseteq}.
                    $f(n)\in\reals$ by \obv{funs_type_apply}.
                    $\surrealLt{
                        \metricDistance[\RealMetric](f(x),f(n))}
                        {\surrealOne}$.
                    Hence $\surrealLe{\abs{f(x)}}{K}$ by
                        \auto{real_unit_close_to_bounded_is_bounded}.
                    Follows by \auto.
            \end{byCase}
        \end{subproof}
    \end{subproof}
    Follows by \auto{real_bounded}.
\end{proof}

\begin{corollary}\label{psi_unbounded_continuous_map_unbounded_on_naturals}
    Suppose $P$ is an almost-disjoint family on $\naturals$.
    Suppose $f$ is continuous from $\psiSpace{P}$ to $\RealTopology$.
    Suppose $\ran{f}$ is not real bounded.
    Then $\img{f}{\naturals}$ is not real bounded.
\end{corollary}
\begin{proof}
    Suppose $\img{f}{\naturals}$ is real bounded.
    Hence $\ran{f}$ is real bounded by
        \obv{psi_continuous_real_map_bounded_if_naturals_bounded}.
    Contradiction.
\end{proof}

\begin{lemma}\label{large_value_choice_range_is_infinite}
    Let $f\in\funs{X}{\reals}$.
    Suppose $\naturals\subseteq X$.
    Let $s\in\funs{\naturals}{\naturals}$.
    Suppose for every natural number $k$ we have
        $\surrealLt{\ordinalSurreal{k}}{\abs{f(s(k))}}$.
    Then $\ran{s}$ is infinite.
\end{lemma}
\begin{proof}
    $f$ and $s$ are functions and $\dom{f}=X$ and
        $\dom{s}=\naturals$ by \obv{funs_is_function,funs}.
    $\ran{s}\subseteq\naturals$ by \obv{funs_ran}.
    Hence $\ran{s}\subseteq X$ by \obv{subseteq_transitive}.
    Hence $\ran{s}\subseteq\dom{f}$.
    Suppose $\ran{s}$ is finite.
    Hence $\img{f}{\ran{s}}$ is finite by
        \obv{finite_function_image}.
    $\img{f}{\ran{s}}\subseteq\ran{f}$ by
        \obv{img_subseteq_ran}.
    $\ran{f}\subseteq\reals$ by \obv{funs_ran}.
    Hence $\img{f}{\ran{s}}\subseteq\reals$ by
        \obv{subseteq_transitive}.
    Take $N\in\naturals$ such that for every
        $y\in\img{f}{\ran{s}}$ we have
            $\surrealLe{\abs{y}}{\ordinalSurreal{N}}$
        by \obv{finite_real_set_has_natural_bound}.
    $s(N)\in\ran{s}$ by \auto{fun_ran_iff}.
    Hence $s(N)\in\dom{f}\inter\ran{s}$ by \obv{inter,subseteq}.
    Hence $f(s(N))\in\img{f}{\ran{s}}$ by
        \obv{img_of_function_intro}.
    Hence $\surrealLe{\abs{f(s(N))}}{\ordinalSurreal{N}}$.
    $\surrealLt{\ordinalSurreal{N}}{\abs{f(s(N))}}$.
    $f(s(N))\in\reals$ by \obv{funs_type_apply,subseteq}.
    Hence $\surreal{\abs{f(s(N))}}$ and
        $\surreal{\ordinalSurreal{N}}$ by
        \auto{surreal_real_is_surreal,surreal_abs_is_surreal,
            ordinal_surreal_is_surreal,natural_number_is_ordinal}.
    Contradiction by \auto{surreal_le,surreal_lt_asymmetric,
        surreal_lt_irreflexive}.
\end{proof}

\begin{theorem}\label{psi_blocking_implies_pseudocompact}
    Suppose $P$ is an almost-disjoint family on $\naturals$.
    Suppose $P$ blocks the infinite subsets of $\naturals$.
    Then $\psiSpace{P}$ is pseudocompact.
\end{theorem}
\begin{proof}
    $\psiSpace{P}$ is a topological space by
        \obv{psi_space_is_topological_space}.
    Show that for every $f$ such that
        $f$ is continuous from $\psiSpace{P}$ to $\RealTopology$
        we have $\ran{f}$ is real bounded.
    \begin{subproof}
        Fix $f$.
        Assume $f$ is continuous from $\psiSpace{P}$ to $\RealTopology$.
        Suppose not.
        $f\in\funs{\carrier[\psiSpace{P}]}{\carrier[\RealTopology]}$
            by \obv{continuous_map}.
        Hence $f\in\funs{\psiCarrier{P}}{\reals}$ by
            \auto{psi_space_carrier,real_topology_carrier}.
        Hence $f$ is a function and $\dom{f}=\psiCarrier{P}$ by
            \obv{funs_is_function,funs}.
        $\img{f}{\naturals}$ is not real bounded by
            \obv{psi_unbounded_continuous_map_unbounded_on_naturals}.
        $\naturals\subseteq\psiCarrier{P}$ by
            \obv{naturals_subset_psi_carrier}.
        Hence $\naturals\subseteq\dom{f}$.
        Show that for every natural number $k$ there exists
            $n\in\naturals$ such that
                $\surrealLt{\ordinalSurreal{k}}{\abs{f(n)}}$.
        \begin{subproof}
            Follows by
                \obv{function_image_unbounded_iff_large_values}.
        \end{subproof}
        Let $R=\{p\in\naturals\times\naturals\mid
            \surrealLt{\ordinalSurreal{\fst{p}}}{\abs{f(\snd{p})}}\}$.
        $R$ is a relation by \auto{relation,times_elem_is_tuple}.
        Show that for every $k\in\naturals$ there exists $n$ such that
            $k\mathrel{R}n$.
        \begin{subproof}
            Fix $k\in\naturals$.
            Take $n\in\naturals$ such that
                $\surrealLt{\ordinalSurreal{k}}{\abs{f(n)}}$.
            $(k,n)\in\naturals\times\naturals$ by
                \obv{times_tuple_intro}.
            $\fst{(k,n)}=k$ and $\snd{(k,n)}=n$ by
                \obv{fst_eq,snd_eq}.
            Hence $(k,n)\in R$ by \auto.
            Follows by \auto.
        \end{subproof}
        Take $s$ such that $s$ is a function on $\naturals$ and
            for every $k\in\naturals$ we have $k\mathrel{R}s(k)$
            by \obv{axiom_of_choice}.
        Hence $s$ is a function and $\dom{s}=\naturals$.
        Show that for every $k\in\naturals$ we have
            $s(k)\in\naturals$ and
            $\surrealLt{\ordinalSurreal{k}}{\abs{f(s(k))}}$.
        \begin{subproof}
            Fix $k\in\naturals$.
            $(k,s(k))\in R$.
            Hence $(k,s(k))\in\naturals\times\naturals$ and
                $\surrealLt{\ordinalSurreal{\fst{(k,s(k))}}}
                    {\abs{f(\snd{(k,s(k))})}}$.
            $\fst{(k,s(k))}=k$ and $\snd{(k,s(k))}=s(k)$ by
                \obv{fst_eq,snd_eq}.
            Follows by \obv{times_tuple_elim}.
        \end{subproof}
        Hence $s$ is a function to $\naturals$ by \auto.
        Hence $s\in\funs{\naturals}{\naturals}$ by \obv{funs_intro}.
        Let $B=\ran{s}$.
        $B\subseteq\naturals$ by \auto{funs_ran}.
        $B$ is infinite by \auto{large_value_choice_range_is_infinite}.
        Take $A\in P$ such that $A\inter B$ is infinite by
            \obv{blocks_infinite_subsets_of}.
        $A\in\carrier[\psiSpace{P}]$ by
            \auto{psi_family_subset_psi_carrier,psi_space_carrier,subseteq}.
        Hence $f(A)\in\reals$ by
            \auto{funs_type_apply,real_topology_carrier,
                psi_space_carrier}.
        Take $M\in\naturals$ such that for every $y\in\reals$ we have if
            $\surrealLt{\metricDistance[\RealMetric](f(A),y)}
                {\surrealOne}$, then
            $\surrealLe{\abs{y}}{\ordinalSurreal{M}}$ by
            \obv{real_unit_ball_has_natural_absolute_bound}.
        Take $F$ such that $F\subseteq\naturals$ and $F$ is finite and
            for every $n\in A\setminus F$ we have
                $\surrealLt{\metricDistance[\RealMetric](f(A),f(n))}
                    {\surrealOne}$ by
            \obv{psi_continuous_real_map_eventually_unit_close}.
        Show that every element of $A\inter B$ is an element of
            $F\union\img{s}{M}$.
        \begin{subproof}
            Fix $x\in A\inter B$.
            $x\in A$ and $x\in B$ by \obv{inter}.
            \begin{byCase}
                \caseOf{$x\in F$.}
                    Follows by \obv{union_intro_left}.
                \caseOf{$x\notin F$.}
                    Hence $x\in A\setminus F$ by \obv{setminus}.
                    Take $k\in\dom{s}$ such that $s(k)=x$ by
                        \auto{fun_ran_iff}.
                    Hence $k\in\naturals$.
                    $\surrealLt{\ordinalSurreal{k}}{\abs{f(s(k))}}$
                        by \auto.
                    $\surrealLt{
                        \metricDistance[\RealMetric](f(A),f(x))}
                        {\surrealOne}$.
                    $A\subseteq\naturals$ by
                        \obv{almost_disjoint_members_subset}.
                    Hence $x\in\naturals$ by \obv{subseteq}.
                    Hence $x\in\psiCarrier{P}$ by
                        \obv{naturals_subset_psi_carrier,subseteq}.
                    Hence $f(x)\in\reals$ by \obv{funs_type_apply}.
                    Hence $\surrealLe{\abs{f(x)}}{\ordinalSurreal{M}}$.
                    Hence $\surrealLt{\ordinalSurreal{k}}
                        {\ordinalSurreal{M}}$ by
                        \auto{surreal_lt_le_transitive,
                            ordinal_surreal_is_surreal,
                            natural_number_is_ordinal,
                            surreal_real_is_surreal,surreal_abs_is_surreal}.
                    Hence $k\in M$ by
                        \auto{ordinal_surreal_lt_iff,
                            natural_number_is_ordinal}.
                    $M\subseteq\naturals$ by
                        \auto{naturals_is_transitive_set,
                            transitiveset_iff_subseteq}.
                    Hence $k\in\dom{s}\inter M$ by \obv{inter,subseteq}.
                    Hence $s(k)\in\img{s}{M}$ by
                        \obv{img_of_function_intro}.
                    Hence $x\in\img{s}{M}$ by \auto.
                    Follows by \obv{union_intro_right}.
            \end{byCase}
        \end{subproof}
        Hence $A\inter B\subseteq F\union\img{s}{M}$ by
            \obv{subseteq}.
        $M$ is finite by \obv{natural_number_is_finite}.
        $M\subseteq\dom{s}$ by
            \auto{naturals_is_transitive_set,transitiveset_iff_subseteq,
                subseteq}.
        Hence $\img{s}{M}$ is finite by \obv{finite_function_image}.
        Hence $F\union\img{s}{M}$ is finite by
            \obv{union_of_finite_is_finite}.
        Hence $A\inter B$ is finite by
            \obv{subseteq_finite_is_finite}.
        Contradiction.
    \end{subproof}
    Follows by \obv{is_pseudocompact}.
\end{proof}

\begin{theorem}\label{psi_not_blocking_implies_not_pseudocompact}
    Suppose $P$ is an almost-disjoint family on $\naturals$.
    Suppose $P$ does not block the infinite subsets of $\naturals$.
    Then $\psiSpace{P}$ is not pseudocompact.
\end{theorem}
\begin{proof}
    Take $B$ such that $B$ is almost disjoint from $P$ on $\naturals$
        by \obv{not_blocking_iff_almost_disjoint_witness}.
    Hence $B\subseteq\naturals$ and $B$ is infinite and
        for every $A\in P$ we have $A\inter B$ is finite
        by \obv{almost_disjoint_from_family_on}.
    Hence $B$ is countably infinite by
        \obv{infinite_subset_of_naturals_is_countably_infinite}.
    Take $e$ such that $e\in\bijections{\naturals}{B}$ by
        \auto{countably_infinite_iff_natural_bijection}.
    $e\in\funs{\naturals}{B}$ and $e$ is injective by
        \auto{bijections,surj}.
    Hence $e$ is a function and $\dom{e}=\naturals$ by
        \obv{funs_is_function,funs}.
    For all $m,n\in\naturals$ such that $e(m)=e(n)$ we have $m=n$ by
        \auto{injective_function}.
    Hence $e\in\inj{\naturals}{B}$ by \obv{inj}.
    $B\subseteq\psiCarrier{P}$ by
        \auto{naturals_subset_psi_carrier,subseteq_transitive}.
    Hence $e\in\inj{\naturals}{\psiCarrier{P}}$ by
        \obv{inj_weaken_codom}.
    $\naturalRealEmbedding\in\funs{\naturals}{\reals}$ by
        \obv{natural_real_embedding_type}.
    $\surrealZero\in\reals$ by \obv{surreal_zero_is_real}.
    Let $g=\functionDomainExtension{e}{\naturalRealEmbedding}
        {\psiCarrier{P}}{\surrealZero}$.
    $g\in\funs{\psiCarrier{P}}{\reals}$ by
        \auto{function_domain_extension_type}.
    Hence $g$ is a function and $\dom{g}=\psiCarrier{P}$ by
        \auto{funs_is_function,funs}.
    $\psiSpace{P}$ is a topological space by
        \obv{psi_space_is_topological_space}.
    $\RealTopology$ is a topological space by
        \obv{real_topology_is_topological_space}.
    $\carrier[\psiSpace{P}]=\psiCarrier{P}$ by
        \obv{psi_space_carrier}.
    $\carrier[\RealTopology]=\reals$ by \obv{real_topology_carrier}.
    Hence $g\in\funs{\carrier[\psiSpace{P}]}
        {\carrier[\RealTopology]}$.
    Show that for every $x\in\carrier[\psiSpace{P}]$ there exists
        $U\in\opens[\psiSpace{P}]$ such that $x\in U$ and
        for every $y\in U$ we have $g(y)=g(x)$.
    \begin{subproof}
        Fix $x\in\carrier[\psiSpace{P}]$.
        Hence $x\in\psiCarrier{P}$ by \obv{psi_space_carrier}.
        Hence $x\in\naturals$ or $x\in P$ by
            \obv{psi_carrier_point_cases}.
        Show that there exists $U\in\opens[\psiSpace{P}]$ such that
            $x\in U$ and for every $y\in U$ we have $g(y)=g(x)$.
        \begin{subproof}
            \begin{byCase}
                \caseOf{$x\in\naturals$.}
                    $\{x\}$ is open in $\psiSpace{P}$ by
                        \obv{psi_natural_point_isolated}.
                    Hence $\{x\}\in\opens[\psiSpace{P}]$ by
                        \obv{is_open}.
                    $x\in\{x\}$ by \obv{singleton_intro}.
                    For every $y\in\{x\}$ we have $g(y)=g(x)$ by
                        \auto{singleton_iff}.
                    Follows by \auto.
                \caseOf{$x\in P$.}
                    Let $F=x\inter B$.
                    $F$ is finite.
                    $x\subseteq\naturals$ by
                        \obv{almost_disjoint_members_subset}.
                    $x$ is infinite by
                        \obv{almost_disjoint_members_infinite}.
                    Hence $F\subseteq\naturals$ by \auto{inter,subseteq}.
                    $\psiNeighbourhood{x}{F}\in\psiBasis{P}$ by
                        \obv{psi_neighbourhood_in_basis}.
                    $\psiBasis{P}$ is an open basis of $\psiSpace{P}$ by
                        \obv{psi_basis_is_open_basis}.
                    Hence $\psiBasis{P}\subseteq\opens[\psiSpace{P}]$ by
                        \obv{open_basis}.
                    Hence $\psiNeighbourhood{x}{F}
                        \in\opens[\psiSpace{P}]$ by \obv{subseteq}.
                    $x\in\psiNeighbourhood{x}{F}$ by
                        \obv{psi_neighbourhood_contains_center}.
                    $\psiNeighbourhood{x}{x\inter B}\inter B=\emptyset$
                        by \obv{psi_neighbourhood_avoids_natural_set}.
                    Hence $\psiNeighbourhood{x}{F}\inter B=\emptyset$
                        by \auto.
                    $e\in\surj{\naturals}{B}$ by
                        \obv{bijections}.
                    Hence $\ran{e}=B$ by \obv{ran_of_surj}.
                    Show that for every $y\in\psiNeighbourhood{x}{F}$
                        we have $g(y)=\surrealZero$.
                    \begin{subproof}
                        Fix $y\in\psiNeighbourhood{x}{F}$.
                        $y\notin B$ by \auto{inter,emptyset}.
                        $\psiNeighbourhood{x}{F}\subseteq\psiCarrier{P}$
                            by \obv{psi_neighbourhood_subset_carrier}.
                        Hence $y\in\psiCarrier{P}$ by \obv{subseteq}.
                        Hence $y\in\psiCarrier{P}\setminus\ran{e}$ by
                            \auto{setminus}.
                        Follows by
                            \auto{function_domain_extension_off_range}.
                    \end{subproof}
                    Hence $g(x)=\surrealZero$.
                    Hence for every $y\in\psiNeighbourhood{x}{F}$
                        we have $g(y)=g(x)$ by \auto.
                    Follows by \auto.
            \end{byCase}
        \end{subproof}
    \end{subproof}
    Hence $g$ is locally constant from $\psiSpace{P}$ to $\RealTopology$
        by \obv{locally_constant_map}.
    Hence $g$ is continuous from $\psiSpace{P}$ to $\RealTopology$ by
        \obv{locally_constant_map_is_continuous}.
    Show that for every $z\in\ran{\naturalRealEmbedding}$ we have
        $z\in\ran{g}$.
    \begin{subproof}
        Fix $z\in\ran{\naturalRealEmbedding}$.
        $\naturalRealEmbedding$ is a function by
            \auto{natural_real_embedding_type,funs_is_function}.
        Take $k\in\dom{\naturalRealEmbedding}$ such that
            $\apply{\naturalRealEmbedding}{k}=z$ by \auto{fun_ran_iff}.
        Hence $k\in\naturals$ by
            \auto{natural_real_embedding_type,funs}.
        $\apply{g}{e(k)}=\apply{\naturalRealEmbedding}{k}$ by
            \auto{function_domain_extension_on_range}.
        $e(k)\in\psiCarrier{P}$ by
            \obv{inj,funs_type_apply}.
        Hence $e(k)\in\dom{g}$.
        Hence $\apply{g}{e(k)}\in\ran{g}$ by \auto{fun_ran_iff}.
    \end{subproof}
    Hence $\ran{\naturalRealEmbedding}\subseteq\ran{g}$ by
        \obv{subseteq}.
    Show that $\ran{g}$ is not real bounded.
    \begin{subproof}
        Suppose $\ran{g}$ is real bounded.
        Hence $\ran{\naturalRealEmbedding}$ is real bounded by
            \obv{subset_of_real_bounded_is_bounded}.
        $\ran{\naturalRealEmbedding}$ is not real bounded by
            \obv{natural_real_embedding_range_unbounded}.
        Contradiction.
    \end{subproof}
    Follows by
        \auto{not_pseudocompact_iff_unbounded_continuous_real_map}.
\end{proof}

\begin{theorem}\label{psi_space_pseudocompact_iff_blocking}
    Suppose $P$ is an almost-disjoint family on $\naturals$.
    Then $\psiSpace{P}$ is pseudocompact iff
        $P$ blocks the infinite subsets of $\naturals$.
\end{theorem}
\begin{proof}
    Show that if $\psiSpace{P}$ is pseudocompact, then
        $P$ blocks the infinite subsets of $\naturals$.
    \begin{subproof}
        Assume $\psiSpace{P}$ is pseudocompact.
        Suppose not.
        Hence $\psiSpace{P}$ is not pseudocompact by
            \obv{psi_not_blocking_implies_not_pseudocompact}.
        Contradiction.
    \end{subproof}
    Show that if $P$ blocks the infinite subsets of $\naturals$, then
        $\psiSpace{P}$ is pseudocompact.
    \begin{subproof}
        Follows by \obv{psi_blocking_implies_pseudocompact}.
    \end{subproof}
    Follows by \auto.
\end{proof}

\begin{theorem}\label{psi_space_pseudocompact_iff_maximal_family}
    Suppose $P$ is an almost-disjoint family on $\naturals$.
    Then $\psiSpace{P}$ is pseudocompact iff
        $P$ is a maximal almost-disjoint family on $\naturals$.
\end{theorem}
\begin{proof}
    $\psiSpace{P}$ is pseudocompact iff
        $P$ blocks the infinite subsets of $\naturals$ by
        \obv{psi_space_pseudocompact_iff_blocking}.
    $P$ is a maximal almost-disjoint family on $\naturals$ iff
        $P$ blocks the infinite subsets of $\naturals$ by
        \obv{maximal_almost_disjoint_iff_blocking}.
\end{proof}

\begin{theorem}\label{psi_space_pseudocompact_iff_infinite_maximal_family}
    Suppose $P$ is an infinite almost-disjoint family on $\naturals$.
    Then $P$ is an infinite maximal almost-disjoint family on $\naturals$ iff
        $\psiSpace{P}$ is pseudocompact.
\end{theorem}
\begin{proof}
    $P$ is an almost-disjoint family on $\naturals$.
    $P$ is an infinite maximal almost-disjoint family on $\naturals$ iff
        $P$ blocks the infinite subsets of $\naturals$ by
        \obv{mad_iff_blocking}.
    $\psiSpace{P}$ is pseudocompact iff
        $P$ blocks the infinite subsets of $\naturals$ by
        \obv{psi_space_pseudocompact_iff_blocking}.
\end{proof}

\begin{theorem}\label{psi_space_structural_theorem}
    Suppose $P$ is an almost-disjoint family on $\naturals$.
    Then $\psiSpace{P}$ is Hausdorff and
        $\psiSpace{P}$ has a clopen basis and
        $\psiSpace{P}$ is completely regular and
        $\psiSpace{P}$ is Tychonoff and
        $\psiSpace{P}$ is locally compact and
        $\psiSpace{P}$ is separable and
        $\psiSpace{P}$ is first-countable.
\end{theorem}
\begin{proof}
    Follows by
        \obv{psi_space_is_hausdorff,psi_space_has_clopen_basis,
            psi_space_is_completely_regular,psi_space_is_tychonoff,
            psi_space_is_locally_compact,psi_space_is_separable,
            psi_space_is_first_countable}.
\end{proof}

\begin{theorem}\label{isbell_mrowka_space_theorem}
    Let $P$ be an infinite almost-disjoint family on $\naturals$.
    Then $\psiSpace{P}$ is Hausdorff and
        $\psiSpace{P}$ has a clopen basis and
        $\psiSpace{P}$ is Tychonoff and
        $\psiSpace{P}$ is locally compact and
        $\psiSpace{P}$ is separable and
        $\psiSpace{P}$ is first-countable and
        $\psiSpace{P}$ is not countably compact and
        ($P$ is an infinite maximal almost-disjoint family on $\naturals$ iff
            $\psiSpace{P}$ is pseudocompact).
\end{theorem}
\begin{proof}
    $P$ is an almost-disjoint family on $\naturals$.
    Follows by
        \obv{psi_space_is_hausdorff,psi_space_has_clopen_basis,
            psi_space_is_tychonoff,psi_space_is_locally_compact,
            psi_space_is_separable,psi_space_is_first_countable,
            infinite_almost_disjoint_psi_space_is_not_countably_compact,
            psi_space_pseudocompact_iff_infinite_maximal_family}.
\end{proof}

\begin{corollary}\label{isbell_mrowka_space_is_not_compact}
    Let $P$ be an infinite almost-disjoint family on $\naturals$.
    Then $\psiSpace{P}$ is not compact.
\end{corollary}
\begin{proof}
    Follows by \obv{infinite_almost_disjoint_psi_space_is_not_compact}.
\end{proof}

\begin{theorem}\label{maximal_isbell_mrowka_space_exists}
    There exists an infinite almost-disjoint family $M$ on $\naturals$ such that
        $M$ is an infinite maximal almost-disjoint family on $\naturals$ and
        $\psiSpace{M}$ is Tychonoff and
        $\psiSpace{M}$ is locally compact and
        $\psiSpace{M}$ has a clopen basis and
        $\psiSpace{M}$ is separable and
        $\psiSpace{M}$ is first-countable and
        $\psiSpace{M}$ is pseudocompact and
        $\psiSpace{M}$ is not countably compact.
\end{theorem}
\begin{proof}
    Take $M$ such that
        $M$ is an infinite maximal almost-disjoint family on $\naturals$ by
        \obv{infinite_maximal_almost_disjoint_family_on_naturals_exists}.
    Hence $M$ is an infinite almost-disjoint family on $\naturals$ by
        \obv{maximal_almost_disjoint_family_on}.
    Hence $M$ is an almost-disjoint family on $\naturals$.
    $\psiSpace{M}$ is Tychonoff by \obv{psi_space_is_tychonoff}.
    $\psiSpace{M}$ is locally compact by
        \obv{psi_space_is_locally_compact}.
    $\psiSpace{M}$ has a clopen basis by
        \obv{psi_space_has_clopen_basis}.
    $\psiSpace{M}$ is separable by \obv{psi_space_is_separable}.
    $\psiSpace{M}$ is first-countable by
        \obv{psi_space_is_first_countable}.
    $\psiSpace{M}$ is pseudocompact by
        \obv{psi_space_pseudocompact_iff_infinite_maximal_family}.
    $\psiSpace{M}$ is not countably compact by
        \obv{infinite_almost_disjoint_psi_space_is_not_countably_compact}.
    Follows by \auto.
\end{proof}

\begin{theorem}\label{isbell_mrowka_example_exists}
    There exists $X$ such that $X$ is a topological space and
        (if $X$ is a topological space, then
            $X$ is Hausdorff and $X$ has a clopen basis and
            $X$ is Tychonoff and
            $X$ is locally compact and $X$ is separable and
            $X$ is first-countable and $X$ is pseudocompact and
            $X$ is not countably compact).
\end{theorem}
\begin{proof}
    Take $M$ such that
        $M$ is an infinite maximal almost-disjoint family on $\naturals$ by
        \obv{infinite_maximal_almost_disjoint_family_on_naturals_exists}.
    $M$ is an infinite almost-disjoint family on $\naturals$ by
        \obv{maximal_almost_disjoint_family_on}.
    $M$ is an almost-disjoint family on $\naturals$.
    $\psiSpace{M}$ is a topological space by
        \obv{psi_space_is_topological_space}.
    $\psiSpace{M}$ is Hausdorff by \obv{psi_space_is_hausdorff}.
    $\psiSpace{M}$ has a clopen basis by
        \obv{psi_space_has_clopen_basis}.
    $\psiSpace{M}$ is Tychonoff by \obv{psi_space_is_tychonoff}.
    $\psiSpace{M}$ is locally compact by
        \obv{psi_space_is_locally_compact}.
    $\psiSpace{M}$ is separable by \obv{psi_space_is_separable}.
    $\psiSpace{M}$ is first-countable by
        \obv{psi_space_is_first_countable}.
    $\psiSpace{M}$ is pseudocompact by
        \obv{psi_space_pseudocompact_iff_infinite_maximal_family}.
    $\psiSpace{M}$ is not countably compact by
        \obv{infinite_almost_disjoint_psi_space_is_not_countably_compact}.
    Show that if $\psiSpace{M}$ is a topological space, then
        $\psiSpace{M}$ is Hausdorff and
        $\psiSpace{M}$ has a clopen basis and
        $\psiSpace{M}$ is Tychonoff and
        $\psiSpace{M}$ is locally compact and
        $\psiSpace{M}$ is separable and
        $\psiSpace{M}$ is first-countable and
        $\psiSpace{M}$ is pseudocompact and
        $\psiSpace{M}$ is not countably compact.
    \begin{subproof}
        Assume $\psiSpace{M}$ is a topological space.
    \end{subproof}
    Follows by \auto.
\end{proof}

\begin{theorem}\label{empty_psi_space_is_topological_space}
    $\psiSpace{\emptyset}$ is a topological space.
\end{theorem}
\begin{proof}
    $\emptyset$ is an almost-disjoint family on $\naturals$ by
        \obv{empty_is_almost_disjoint_family}.
    Follows by \obv{psi_space_is_topological_space}.
\end{proof}

\begin{proposition}\label{empty_psi_space_carrier}
    Suppose $\psiSpace{\emptyset}$ is a topological space.
    Then $\carrier[\psiSpace{\emptyset}]=\naturals$.
\end{proposition}
\begin{proof}
    $\carrier[\psiSpace{\emptyset}]=\psiCarrier{\emptyset}$ by
        \obv{psi_space_carrier}.
    $\psiCarrier{\emptyset}=\naturals\union\emptyset$ by
        \obv{psi_carrier}.
    $\naturals\union\emptyset=\naturals$ by \obv{union_emptyset}.
\end{proof}

\begin{proposition}\label{empty_psi_every_natural_singleton_open}
    Suppose $\emptyset$ is an almost-disjoint family on
        $\naturals$.
    Let $n\in\naturals$.
    Then $\{n\}$ is open in $\psiSpace{\emptyset}$.
\end{proposition}
\begin{proof}
    $\emptyset$ is an almost-disjoint family on $\naturals$ by
        \obv{empty_is_almost_disjoint_family}.
    Follows by \obv{psi_natural_point_isolated}.
\end{proof}

\begin{theorem}\label{empty_psi_space_is_discrete}
    Suppose $\emptyset$ is an almost-disjoint family on
        $\naturals$.
    Then $\psiSpace{\emptyset}$ is discrete.
\end{theorem}
\begin{proof}
    $\psiSpace{\emptyset}$ is a topological space by
        \obv{empty_psi_space_is_topological_space}.
    Show that for every $x\in\carrier[\psiSpace{\emptyset}]$ we have
        $\{x\}\in\opens[\psiSpace{\emptyset}]$.
    \begin{subproof}
        Fix $x\in\carrier[\psiSpace{\emptyset}]$.
        Hence $x\in\naturals$ by \obv{empty_psi_space_carrier}.
        $\{x\}$ is open in $\psiSpace{\emptyset}$ by
            \obv{empty_psi_every_natural_singleton_open}.
        Follows by \obv{is_open}.
    \end{subproof}
    Follows by \obv{discrete_iff_every_singleton_open}.
\end{proof}

\begin{theorem}\label{empty_psi_space_is_not_pseudocompact}
    Suppose $\emptyset$ is an almost-disjoint family on
        $\naturals$.
    Then $\psiSpace{\emptyset}$ is not pseudocompact.
\end{theorem}
\begin{proof}
    $\emptyset$ is an almost-disjoint family on $\naturals$ by
        \obv{empty_is_almost_disjoint_family}.
    Show that $\naturals$ is almost disjoint from $\emptyset$ on
        $\naturals$.
    \begin{subproof}
        $\naturals\subseteq\naturals$ by \obv{subseteq_refl}.
        $\naturals$ is infinite by \obv{naturals_is_infinite}.
        For every $A\in\emptyset$ we have $A\inter\naturals$ is finite
            by \auto{emptyset}.
        Follows by \obv{almost_disjoint_from_family_on}.
    \end{subproof}
    Hence $\emptyset$ does not block the infinite subsets of $\naturals$
        by \auto{not_blocking_iff_almost_disjoint_witness}.
    Follows by \obv{psi_not_blocking_implies_not_pseudocompact}.
\end{proof}

\begin{theorem}\label{empty_psi_space_is_not_countably_compact}
    Suppose $\emptyset$ is an almost-disjoint family on
        $\naturals$.
    Then $\psiSpace{\emptyset}$ is not countably compact.
\end{theorem}
\begin{proof}
    $\psiSpace{\emptyset}$ is a topological space by
        \obv{empty_psi_space_is_topological_space}.
    $\naturals$ is countable by \obv{naturals_is_countable}.
    $\naturals$ is infinite by \obv{naturals_is_infinite}.
    Hence $\naturals$ is countably infinite by
        \obv{countably_infinite}.
    $\naturals=\carrier[\psiSpace{\emptyset}]$ by
        \obv{empty_psi_space_carrier}.
    Hence $\naturals$ is closed in $\psiSpace{\emptyset}$ by
        \auto{carrier_is_closed}.
    Show that for every $x\in\naturals$ there exists
        $U\in\opens[\psiSpace{\emptyset}]$ such that
            $U\inter\naturals=\{x\}$.
    \begin{subproof}
        Fix $x\in\naturals$.
        $\{x\}$ is open in $\psiSpace{\emptyset}$ by
            \obv{empty_psi_every_natural_singleton_open}.
        Hence $\{x\}\in\opens[\psiSpace{\emptyset}]$ by
            \obv{is_open}.
        $\{x\}\subseteq\naturals$ by \obv{singleton_subset_intro}.
        Hence $\{x\}\inter\naturals=\{x\}$ by
            \obv{inter_absorb_supseteq_right}.
        Follows by \auto.
    \end{subproof}
    $\naturals\subseteq\naturals$ by \obv{subseteq_refl}.
    Hence $\naturals\subseteq\carrier[\psiSpace{\emptyset}]$.
    Hence $\naturals$ is discrete in $\psiSpace{\emptyset}$ by
        \auto{discrete_in_iff_ambient_open_singletons}.
    Follows by
        \obv{countably_infinite_closed_discrete_obstructs_countable_compactness}.
\end{proof}

\begin{proposition}\label{empty_almost_disjoint_family_is_not_maximal}
    $\emptyset$ is not a maximal almost-disjoint family on
        $\naturals$.
\end{proposition}
\begin{proof}
    $\emptyset$ is an almost-disjoint family on $\naturals$ by
        \obv{empty_is_almost_disjoint_family}.
    Show that $\naturals$ is almost disjoint from $\emptyset$ on
        $\naturals$.
    \begin{subproof}
        $\naturals\subseteq\naturals$ by \obv{subseteq_refl}.
        $\naturals$ is infinite by \obv{naturals_is_infinite}.
        For every $A\in\emptyset$ we have $A\inter\naturals$ is finite
            by \auto{emptyset}.
        Follows by \obv{almost_disjoint_from_family_on}.
    \end{subproof}
    Hence $\emptyset$ does not block the infinite subsets of $\naturals$
        by \auto{not_blocking_iff_almost_disjoint_witness}.
    Suppose $\emptyset$ is a maximal almost-disjoint family on
        $\naturals$.
    Hence $\emptyset$ blocks the infinite subsets of $\naturals$ by
        \obv{maximal_almost_disjoint_iff_blocking}.
    Contradiction.
\end{proof}

\begin{proposition}\label{singleton_naturals_is_almost_disjoint_family}
    $\{\naturals\}$ is an almost-disjoint family on $\naturals$.
\end{proposition}
\begin{proof}
    $\naturals\subseteq\naturals$ by \obv{subseteq_refl}.
    Hence $\naturals\in\pow{\naturals}$ by \auto{pow_iff}.
    Hence $\{\naturals\}\subseteq\pow{\naturals}$ by
        \obv{singleton_subset_intro}.
    For every $A\in\{\naturals\}$ we have $A$ is infinite by
        \auto{singleton_iff,naturals_is_infinite}.
    For all $A,B\in\{\naturals\}$ such that $A\neq B$ we have
        $A\inter B$ is finite by \auto{singleton_iff}.
    Follows by \obv{almost_disjoint_family_on}.
\end{proof}

\begin{theorem}\label{singleton_naturals_psi_space_is_topological_space}
    $\psiSpace{\{\naturals\}}$ is a topological space.
\end{theorem}
\begin{proof}
    Follows by
        \obv{singleton_naturals_is_almost_disjoint_family,
            psi_space_is_topological_space}.
\end{proof}

\begin{proposition}\label{singleton_naturals_blocks_infinite_subsets}
    $\{\naturals\}$ blocks the infinite subsets of $\naturals$.
\end{proposition}
\begin{proof}
    Show that for every $B\subseteq\naturals$ such that $B$ is infinite
        there exists $A\in\{\naturals\}$ such that
            $A\inter B$ is infinite.
    \begin{subproof}
        Fix $B\subseteq\naturals$.
        Assume $B$ is infinite.
        $\naturals\in\{\naturals\}$ by \obv{singleton_intro}.
        $\naturals\inter B=B$ by
            \obv{inter_absorb_supseteq_left}.
        Hence $\naturals\inter B$ is infinite.
        Follows by \auto.
    \end{subproof}
    Follows by \obv{blocks_infinite_subsets_of}.
\end{proof}

\begin{proposition}\label{singleton_naturals_is_maximal_almost_disjoint_family}
    $\{\naturals\}$ is a maximal almost-disjoint family on
        $\naturals$.
\end{proposition}
\begin{proof}
    Follows by
        \obv{singleton_naturals_is_almost_disjoint_family,
            singleton_naturals_blocks_infinite_subsets,
            maximal_almost_disjoint_iff_blocking}.
\end{proof}

\begin{proposition}\label{singleton_naturals_is_not_infinite_almost_disjoint_family}
    $\{\naturals\}$ is not an infinite almost-disjoint family on $\naturals$.
\end{proposition}
\begin{proof}
    $\{\naturals\}$ is finite by \obv{singleton_is_finite}.
    Suppose $\{\naturals\}$ is an infinite almost-disjoint family on $\naturals$.
    Hence $\{\naturals\}$ is infinite.
    Contradiction.
\end{proof}

\begin{theorem}\label{singleton_naturals_psi_space_is_compact}
    Suppose $\{\naturals\}$ is an almost-disjoint family on
        $\naturals$.
    Then $\psiSpace{\{\naturals\}}$ is compact.
\end{theorem}
\begin{proof}
    Let $X=\psiSpace{\{\naturals\}}$.
    $X$ is a topological space by
        \obv{singleton_naturals_psi_space_is_topological_space}.
    $\naturals\in\{\naturals\}$ by \obv{singleton_intro}.
    Hence $\naturals\in\psiCarrier{\{\naturals\}}$ by
        \auto{psi_family_subset_psi_carrier,subseteq}.
    Hence $\naturals\in\carrier[X]$ by \auto{psi_space_carrier}.
    Show that for every $U\in\opens[X]$ such that $\naturals\in U$ we have
        $\carrier[X]\setminus U$ is finite.
    \begin{subproof}
        Fix $U\in\opens[X]$.
        Assume $\naturals\in U$.
        $U$ is open in $X$ by \obv{is_open}.
        Hence $\naturals\in\{\naturals\}\inter U$ by \obv{inter}.
        Take $F$ such that $F\subseteq\naturals$ and $F$ is finite and
            $\psiNeighbourhood{\naturals}{F}\subseteq U$ by
            \obv{psi_family_point_open_refinement}.
        Show that every element of $\carrier[X]\setminus U$ is an
            element of $F$.
        \begin{subproof}
            Fix $x\in\carrier[X]\setminus U$.
            $x\in\carrier[X]$ and $x\notin U$ by \obv{setminus}.
            Hence $x\in\psiCarrier{\{\naturals\}}$ by
                \auto{psi_space_carrier}.
            Hence $x\in\naturals$ or $x=\naturals$ by
                \auto{psi_carrier_iff,singleton_iff}.
            $x\notin\psiNeighbourhood{\naturals}{F}$ by
                \auto{subseteq}.
            $x\neq\naturals$.
            Hence $x\in\naturals$.
            Hence $x\in F$ by \auto{psi_neighbourhood,setminus,union}.
        \end{subproof}
        Hence $\carrier[X]\setminus U\subseteq F$ by \obv{subseteq}.
        Follows by \obv{subseteq_finite_is_finite}.
    \end{subproof}
    Follows by \auto{cofinite_neighbourhood_point_implies_compact}.
\end{proof}

\begin{corollary}\label{singleton_naturals_psi_space_is_pseudocompact}
    Suppose $\{\naturals\}$ is an almost-disjoint family on
        $\naturals$.
    Then $\psiSpace{\{\naturals\}}$ is pseudocompact.
\end{corollary}
\begin{proof}
    Follows by
        \obv{singleton_naturals_is_almost_disjoint_family,
            singleton_naturals_is_maximal_almost_disjoint_family,
            psi_space_pseudocompact_iff_maximal_family}.
\end{proof}

\begin{proposition}\label{psi_neighbourhood_finite_set_normalization_regression}
    Suppose $P$ is an almost-disjoint family on $\naturals$.
    Suppose $A\in P$.
    Suppose $F\subseteq\naturals$ and $F$ is finite.
    Then $\psiNeighbourhood{A}{F}
        =\psiNeighbourhood{A}{F\inter A}$.
\end{proposition}
\begin{proof}
    Follows by \obv{psi_neighbourhood_normalize_finite_set}.
\end{proof}
